\subsection{Ingoing Eddington--Finkelstein Einstein--Vlasov Weighted
Particle-in-Cell (WPIC) Module}
\label{sec:ingoing-eddington-finkelstein-wpic}

\[
 \boxed{\dd s^2=-a(v,r)b(v,r)^2\dd v^2
          +2b(v,r)\dd v\,\dd r+r^2\dd\Omega^2.}          \tag{IEF2}
\]
This is the ingoing Eddington--Finkelstein (EF) line element used by VENUS.
The symbol \(\dd s^2\) is the squared spacetime interval; the metric signature
is \((-+++)\).  The coordinate \(v\) is advanced, or ingoing-null, time.
The coordinate \(r\) is areal radius, so a sphere at \(r\) has area
\(4\pi r^2\).  The angles are \(\theta\) and \(\phi\), and
\(\dd\Omega^2=\dd\theta^2+\sin^2\theta\,\dd\phi^2\) is the unit-sphere line
element.  The dimensionless functions \(a(v,r)\) and \(b(v,r)\) determine the
geometry; VENUS uses the future-directed branch \(b>0\).  Geometrized units
\(G=c=1\) are used throughout.

The derivation first establishes the null geometry, then obtains the field
constraints and particle characteristics.  The implementation sections follow
the same order, from source deposition to the metric solve and particle
update.

\subsubsection{The geometry before the equations}
\label{sec:ef-geometry-first}

In the two coordinates \((v,r)\), the metric and its inverse are
\[
 (g_{AB})=
 \begin{pmatrix}-ab^2&b\\ b&0\end{pmatrix},\qquad
 (g^{AB})=
 \begin{pmatrix}0&b^{-1}\\ b^{-1}&a\end{pmatrix}.          \tag{IEF3}
\]
Upper-case indices \(A,B\) run over \(v,r\); lower-case indices \(a,b\) run
over all four spacetime coordinates.  The determinant and invariant
four-volume are
\[
 g:=\det(g_{ab})=-b^2r^4\sin^2\theta,\qquad
 \sqrt{-g}\,\dd v\,\dd r\,\dd\theta\,\dd\phi
 =br^2\sin\theta\,\dd v\,\dd r\,\dd\theta\,\dd\phi.       \tag{IEF4}
\]
The off-diagonal term \(2b\,\dd v\,\dd r\) keeps the \((v,r)\) block
nondegenerate even when \(a=0\).  This simple algebra is the local reason the
chart can remain regular at an apparent horizon.

The Hawking--Misner--Sharp mass \(m(v,r)\), meaning the gravitational mass
enclosed by a symmetry sphere, is defined by
\[
 a=g^{AB}\partial_A r\,\partial_B r
   =1-\frac{2m(v,r)}{r}.                                  \tag{IEF5}
\]
Here \(\partial_A\) is a coordinate partial derivative.  A subscript on a
field will mean the same operation, for example
\(b_r=(\partial b/\partial r)_v\) and
\(a_v=(\partial a/\partial v)_r\).

\paragraph{Immediate comparison with the Schwarzschild chart.}
The companion Schwarzschild-coordinate (SC) module uses the Schwarzschild
line element
\[
 \dd s^2=-e^{2\mu(\tau,r)}\dd\tau^2
          +e^{2\lambda(\tau,r)}\dd r^2+r^2\dd\Omega^2.    \tag{IEF38}
\]
The name Schwarzschild coordinates describes this diagonal spherical gauge;
it does not assume that the spacetime is the static vacuum Schwarzschild
solution.  The coordinate \(\tau\) labels spacelike Schwarzschild slices,
\(\mu(\tau,r)\) is the lapse potential, and \(\lambda(\tau,r)\) is the radial
metric potential.  The two descriptions keep the same areal radius and
angles but assign different time labels to the same event:
\[
       v=v(\tau,r),\qquad r_{\rm EF}=r_{\rm SC}.
\]
An SC surface \(\tau=\mathrm{constant}\) is a simultaneous spacelike slice.
An EF surface \(v=\mathrm{constant}\) is instead an inward-moving light
front.  Thus equal printed values of \(v\) and \(\tau\) do not identify the
same events.

\paragraph{What advanced time and radial null curves mean.}
Set \(\dd\Omega=0\) and \(\dd s^2=0\) in (IEF2).  The radial-null condition
factors as
\[
 \dd v\,(-ab^2\dd v+2b\,\dd r)=0.
\]
It has two future radial families:
\[
 \boxed{\dd v=0\quad\hbox{(ingoing)}},\qquad
 \boxed{\frac{\dd r}{\dd v}=\frac{ab}{2}
        \quad\hbox{(outgoing)}}.                          \tag{IEF6a}
\]
A constant-\(v\) surface is therefore built from ingoing radial light rays.
The word advanced means that this null label is normalized toward the future
at the outer boundary: later ingoing light fronts have larger \(v\).

For later geometric diagnostics choose future null vectors
\[
 k^a=(\partial_v)^a+\frac{ab}{2}(\partial_r)^a,\qquad
 n^a=-\frac1b(\partial_r)^a,\qquad k^a n_a=-1.            \tag{IEF6}
\]
The operator \(\partial_v\) differentiates while holding \(r,\theta,\phi\)
fixed, and similarly for \(\partial_r\).  Their spherical expansions, meaning
the fractional rates at which the areas of neighboring light fronts change,
are
\[
 \theta_+=\frac{2k(r)}r=\frac{ab}{r},\qquad
 \theta_-=\frac{2n(r)}r=-\frac{2}{br},\qquad
 \theta_+\theta_-=-\frac{2a}{r^2}.                        \tag{IEF7}
\]
Consequently an outer marginally trapped sphere satisfies
\[
 a=0\quad\Longleftrightarrow\quad 2m/r=1.                 \tag{IEF8}
\]
At this apparent horizon an outgoing radial light ray has
\(\dd r/\dd v=0\).  In a trapped region \(a<0\), both future radial-null
directions move toward smaller areal radius.  An apparent horizon is a
slice-local diagnostic; the event horizon is a global causal boundary and
need not coincide with it.

\subsubsection{Eikonal construction: from static Schwarzschild to a dynamical map}
\label{sec:ef-eikonal-construction}

Before specializing to any coordinate, let \(\Phi(x^a)\) be a scalar whose
level sets are proposed null hypersurfaces.  Its normal covector is
\(k_a=\nabla_a\Phi=\partial_a\Phi\), where \(\nabla_a\) is the spacetime
covariant derivative and the second equality holds because \(\Phi\) is a
scalar.  The general eikonal equation is
\[
 \boxed{g^{ab}\partial_a\Phi\,\partial_b\Phi=0.}           \tag{IEF-EIK}
\]
It says that the normal has zero squared length.  A null hypersurface is
special because its normal is also tangent to the light rays that generate
the surface.

\paragraph{Static Schwarzschild derivation.}
For a vacuum mass \(M>0\), define
\(f_{\rm S}(r):=1-2M/r\).  The exterior Schwarzschild metric is
\[
 \dd s^2=-f_{\rm S}\dd\tau^2+f_{\rm S}^{-1}\dd r^2
          +r^2\dd\Omega^2,\qquad r>2M.                    \tag{IEF-SCH}
\]
Radial light rays obey
\[
 0=-f_{\rm S}\dd\tau^2+f_{\rm S}^{-1}\dd r^2
 \quad\Longrightarrow\quad
 \frac{\dd r}{\dd\tau}=\pm f_{\rm S}.                     \tag{IEF-SCH-null}
\]
The minus sign is ingoing and the plus sign is outgoing.  Introduce the
tortoise radius \(r_\star\) by
\[
 \frac{\dd r_\star}{\dd r}=\frac1{f_{\rm S}},\qquad
 r_\star=r+2M\log\left|\frac{r}{2M}-1\right|+\mathrm{constant}. \tag{IEF40b-vac}
\]
The advanced and retarded null labels are
\[
       v=\tau+r_\star,\qquad u_{\rm ret}=\tau-r_\star.
\]
Along an ingoing ray \(\dd r/\dd\tau=-f_{\rm S}\), hence \(\dd v=0\);
along an outgoing ray \(\dd u_{\rm ret}=0\).  Substituting
\(\dd\tau=\dd v-\dd r/f_{\rm S}\) cancels the divergent diagonal
\(\dd r^2/f_{\rm S}\) term and gives
\[
 \dd s^2=-f_{\rm S}\dd v^2+2\dd v\,\dd r+r^2\dd\Omega^2.
\]
Thus the static vacuum case of (IEF2) has
\(a=f_{\rm S}\) and \(b=1\).  The mixed term records the change to a null time coordinate.  Particle
motion follows from the same spacetime metric expressed in these coordinates.

Schwarzschild time fails at \(r=2M\) because
\(g_{rr}=f_{\rm S}^{-1}\) and \(r_\star\) diverge there.  In EF variables
the determinant (IEF4), inverse (IEF3), null slopes (IEF6a), Einstein
constraints below, and massive-particle characteristics remain finite when
\(a=0\), provided \(0<b<\infty\) and the future momentum variable \(P\)
remains positive.  This permits numerical evolution through an apparent
horizon and into a trapped region.  It does not regularize a curvature
singularity, turn an apparent horizon into an event-horizon measurement,
supply data at the true origin, or justify arbitrary removal of the interior.
Post-horizon evolution still requires a causally placed excision surface,
fixed-mass bookkeeping, finite fields, and convergence under changes of
resolution and excision controls.

\paragraph{The dynamical eikonal equation.}
For the general SC metric (IEF38), an advanced label \(v(\tau,r)\) obeys
\[
 0=-e^{-2\mu}v_\tau^2+e^{-2\lambda}v_r^2,\qquad
 v_r=+e^{\lambda-\mu}v_\tau.                              \tag{IEF38a}
\]
The positive root is the ingoing family because a level set satisfies
\[
 \left.\frac{\dd r}{\dd\tau}\right|_v
 =-\frac{v_\tau}{v_r}=-e^{\mu-\lambda}<0.                 \tag{IEF38b}
\]
Define the positive map factor \(A:=v_\tau>0\), the inverse radial light speed
\(h:=e^{\lambda-\mu}>0\), and the radial light speed
\(c:=h^{-1}=e^{\mu-\lambda}>0\).  Then
\[
 \dd v=A(\dd\tau+h\,\dd r),\qquad
 \partial_r A=\partial_\tau(Ah).                          \tag{IEF38c}
\]
The second equation is the integrability condition
\(\partial_r v_\tau=\partial_\tau v_r\).  It is why a dynamical map requires
spacetime history rather than a single metric snapshot.

Equivalently, \(v\) solves the transport equation
\[
 v_\tau-cv_r=0,\qquad
 \frac{\dd r_-}{\dd\tau}=-c(\tau,r_-),\qquad
 \left.\frac{\dd v}{\dd\tau}\right|_{r_-}=0,              \tag{IEF38c-1}
\]
whose characteristics \(r_-(\tau)\) are ingoing radial light rays.  On an
initial SC slice \(\tau=\tau_0\) and a vacuum worldtube \(r=R_B\), a compatible
normalization is
\[
 v(\tau_0,r)=\tau_0+v_0-\int_r^{R_B}
 e^{\lambda(\tau_0,\rho)-\mu(\tau_0,\rho)}\,\dd\rho,
 \qquad
 v(\tau,R_B)=\tau+v_0.                                   \tag{IEF38c-2}
\]
Here \(R_B\) is the fixed map-boundary radius, \(\rho\) is only the dummy
integration radius, and \(v_0\) fixes the additive time origin.  The boundary
condition makes \(A=1\) on the worldtube.  A numerical construction must also
satisfy the usual corner compatibility where the initial slice meets the
worldtube.

The complete Jacobians are
\[
 \frac{\partial(v,r)}{\partial(\tau,r)}
 =\begin{pmatrix}A&Ah\\0&1\end{pmatrix},\qquad
 \frac{\partial(\tau,r)}{\partial(v,r)}
 =\begin{pmatrix}F&-h\\0&1\end{pmatrix},\qquad F:=A^{-1}>0, \tag{IEF38c-5}
\]
Consequently, the derivative operators at the same event satisfy
\[
 \left.\partial_v\right|_r=F\left.\partial_\tau\right|_r,
 \qquad
 \left.\partial_r\right|_v
 =\left.\partial_r\right|_\tau-h\left.\partial_\tau\right|_r. \tag{IEF38c-6}
\]
Substitution of \(\dd\tau=F\dd v-h\dd r\) into (IEF38) yields
\[
 \dd s^2=-F^2e^{2\mu}\dd v^2
       +2F e^{\lambda+\mu}\dd v\,\dd r+r^2\dd\Omega^2,
 \qquad
 a=e^{-2\lambda},\quad b=F e^{\lambda+\mu}.               \tag{IEF38d}
\]
Thus \(a\) is the same invariant radial-mass combination in both charts,
whereas \(b\) contains both SC metric potentials and the dynamical
normalization \(F\).  The upper-case \(A=v_\tau\) is a map Jacobian and must
not be confused with the lower-case EF field \(a=1-2m/r\).

In a static geometry \(A=F=1\), giving
\[
 v=\tau+\int^r e^{\lambda-\mu}\dd r'
   =\tau+\int^r\frac{\dd r'}{a(r')b(r')}.                 \tag{IEF39}
\]
The static metric functions are therefore
\[
 a=e^{-2\lambda},\qquad b=e^{\lambda+\mu}.               \tag{IEF40}
\]
Conversely, a local diagonal polar time can be reconstructed from EF data only
where \(a>0\):
\[
 \dd\tau=F(v,r)\left(\dd v-\frac{\dd r}{ab}\right),\qquad
 \partial_r F=-\partial_v\left(\frac{F}{ab}\right).       \tag{IEF40d}
\]
The factor \(1/(ab)\) makes this inverse construction singular at \(a=0\).
Therefore SC--EF comparison is an exterior overlap test; the EF evolution
itself does not use this inverse map after horizon formation.

\subsubsection{Neutral spherical matter and scope}
\label{sec:ef-neutral-scope}

The module evolves a collisionless, spherically symmetric ensemble of
uncharged unit-rest-mass particles.  Collisionless means that particles
interact through the self-consistent metric.  Material collisions are
introduced separately in the relativistic Boltzmann chapter.  Its implementation is
\path{cpu/include/ev/weighted_einstein_vlasov_pic_eddington_finkelstein.hpp}
in the namespace \code{ev::eddington_finkelstein}.  Its marker, metric, and
source data objects are distinct from those of the SC module, so an SC restart
cannot be reinterpreted as EF data.

This is the strict neutral limit of the cited Einstein--Maxwell--Vlasov
formulation:
\[
 \mathfrak e=0,\qquad q_0=0,\qquad Q_{\rm ext}=0,\qquad
 Q_{\rm encl}(v,r)=0,\qquad {\cal F}_{ab}=0.              \tag{IEF1}
\]
Here \(\mathfrak e\) is particle charge, \(q_0\) is charge per unit rest mass,
\(Q_{\rm ext}\) is total exterior charge, \(Q_{\rm encl}\) is enclosed charge,
and \({\cal F}_{ab}\) is the electromagnetic field tensor.  Consequently
there is no Maxwell current or stress, Coulomb mass term, or Lorentz force.
The letter \(F=A^{-1}\) used for the coordinate map is unrelated to the
electromagnetic tensor \({\cal F}_{ab}\).

Every spherical particle carries an angular momentum magnitude \(L\), but its
angular direction is populated isotropically.  The total angular momentum
vector \(\boldsymbol J\) and the component \(T_{v\phi}\) vanish.  A black hole
formed by this module therefore has zero spin; angular momentum threshold
below refers to the particle centrifugal barrier, not Kerr spin.

The continuum equations follow the neutral systems described by East and by
Andr\'easson--Rein.  Cloud shapes, finite-volume quadrature, face-aware
interpolation, fourth-order Runge--Kutta (RK4), midpoint field coupling,
timestep control, and excision thresholds are VENUS numerical choices.  They
must be tested for convergence and are not new continuum assumptions.

\paragraph{Equation and implementation provenance.}
The metric, radial equations, characteristics, and particle sources follow
Appendices A and B of William East,
\href{https://arxiv.org/abs/2511.20567}{arXiv:2511.20567}, in the strict
neutral limit.  The complete neutral characteristic Einstein--Vlasov system
is independently given by H.~Andr\'easson and G.~Rein,
\href{https://arxiv.org/abs/0910.1254}{arXiv:0910.1254}.  East's public \code{evolveEF.c} provides an independent implementation
of these equations. VENUS uses the radial mesh, quadrature and particle
updates derived in this chapter.

\subsubsection{Einstein equations on an ingoing-null slice}
\label{sec:ef-einstein-derivation}

Einstein's equation,
\[
                 G_{ab}=8\pi T_{ab},
\]
equates the Einstein curvature tensor \(G_{ab}\) to the matter stress-energy
tensor \(T_{ab}\).  The mixed \(dv\,dr\) metric term changes coordinate
components, not the underlying tensor.  A null slice also cuts across moving
matter differently from a spacelike SC slice, so its radial mass accumulation
contains both local energy and radial current.

A direct calculation from (IEF2) gives the independent covariant components
\begin{align}
 -\frac{b}{r^2}
 \left(ra_v+rab a_r+a^2b-ab\right)
 &=8\pi T_{vv},                                           \tag{IEF24}\\
 \frac{b}{r^2}\left(ra_r+a-1\right)
 &=8\pi T_{vr},                                           \tag{IEF25}\\
 \frac{2b_r}{rb}
 &=8\pi T_{rr},                                           \tag{IEF26}\\
 \frac{r^2}{2b^2}\left(
 2b_{rv}+b^2a_{rr}+2ab b_{rr}
 -\frac{2b_rb_v}{b}+3b a_rb_r
 +\frac{2b^2a_r}{r}+\frac{2ab b_r}{r}\right)
 &=8\pi T_{\theta\theta}.                                \tag{IEF27}
\end{align}
The \(\phi\phi\) equation is
\(T_{\phi\phi}=\sin^2\theta\,T_{\theta\theta}\).  The first two subscripts on
\(T_{ab}\) name covariant coordinate components, not derivatives.

Equation (IEF26) immediately gives the first radial constraint,
\[
 \boxed{\partial_r\ln b=4\pi rT_{rr}.}                    \tag{IEF28}
\]
Define the positive source combination
\[
 S:=-\frac{2}{b}T_{vr}-aT_{rr}.                           \tag{IEF23a}
\]
Insert this definition into (IEF25), use \(a=1-2m/r\), and combine the result
with (IEF28).  The second radial constraint becomes
\[
 \boxed{\frac1b\partial_r(mb)
       =2\pi r^2(T_{rr}+S).}                              \tag{IEF29}
\]
These two ordinary differential equations are the metric equations solved on
each \(v=\) constant light front.  The pusher needs
\[
 b_r=4\pi rbT_{rr},\qquad
 a_r=\frac{1-a}{r}-4\pi r(aT_{rr}+S).                    \tag{IEF30}
\]
The unused \(vv\) equation supplies the continuum time identity
\[
 a_v=-8\pi r\left(\frac{T_{vv}}b+aT_{vr}\right),\qquad
 m_v=4\pi r^2\left(\frac{T_{vv}}b+aT_{vr}\right).         \tag{IEF31}
\]
In vacuum outside all matter these constraints give the exact exterior
\[
 m=M_{\rm ADM},\qquad b=1,\qquad
 a=1-\frac{2M_{\rm ADM}}r.                               \tag{IEF32}
\]
When the Vlasov equation enforces \(\nabla_aT^{ab}=0\), the contracted Bianchi
identity \(\nabla_aG^{ab}=0\) propagates the remaining Einstein components.
Here \(\nabla_a\) is the metric-compatible covariant derivative.

\paragraph{Equation-by-equation bridge to SC variables.}
Let \(\rho\), \(p_{\rm rad}\), and \(j\) denote energy density, radial pressure,
and outward radial current measured in the local orthonormal frame of the SC
chart.  The SC constraints are
\[
 \left.m_r\right|_\tau=4\pi r^2\rho,\qquad
 \lambda_r=\frac{4\pi r\rho-m/r^2}{a},\qquad
 \mu_r=\frac{m/r^2+4\pi rp_{\rm rad}}{a},\qquad
 \lambda_\tau=-4\pi r e^{\lambda+\mu}j.                  \tag{IEF-SC-E}
\]
At the same event, tensor transformation with (IEF38c-5) gives
\[
 T_{vv}=F^2e^{2\mu}\rho,\qquad
 T_{vr}=-b(\rho+j),\qquad
 T_{rr}=e^{2\lambda}(\rho+2j+p_{\rm rad}),\qquad
 S=\rho-p_{\rm rad}.                                     \tag{IEF45i-2}
\]
The current appears twice in \(T_{rr}\) because the radial basis vector at
fixed \(v\) contains a time-directed part:
\[
 \left.\partial_r\right|_v
 =\left.\partial_r\right|_\tau-h\left.\partial_\tau\right|_r.
\]
Using this derivative and \(b=F e^{\lambda+\mu}\) makes the equivalence
transparent:
\begin{align}
 \left.(\ln b)_r\right|_v
 &=\lambda_r+\mu_r-2h\lambda_\tau
   =4\pi rT_{rr},                                        \tag{IEF-SC-E1}\\
 \left.a_r\right|_v
 &=\frac{1-a}{r}-8\pi r(\rho+j)
   =\frac{1-a}{r}-4\pi r(aT_{rr}+S),                     \tag{IEF-SC-E2}\\
 \left.m_r\right|_v
 &=4\pi r^2(\rho+j)
   =2\pi r^2(aT_{rr}+S),                                 \tag{IEF-SC-E3}\\
 m_v&=F m_\tau=-4\pi r^2F e^{\mu-\lambda}j
   =4\pi r^2\left(\frac{T_{vv}}b+aT_{vr}\right).         \tag{IEF-SC-E4}
\end{align}
Thus (IEF28) is the transformed sum of the SC radial lapse and metric
constraints, (IEF30) is the transformed SC mass constraint, (IEF29) is its
numerically convenient product form, and (IEF31) is the transformed SC flux
equation.  No extra matter source has been introduced by the EF chart.

\paragraph{Radial boundary formulations.}
Let
\[
 Y(r):=m(r)b(r),\qquad
 Y_r=2\pi b r^2(T_{rr}+S).                               \tag{IEF33}
\]
With the asymptotic normalization \(b(\infty)=1\),
\[
 b(r)=\exp\left[-4\pi\int_r^\infty r'T_{rr}(r')\,\dd r'\right]. \tag{IEF34}
\]
The reusable library supports a regular first slice with no point mass:
\[
 Y(0)=0,\qquad
 Y(r)=\int_0^r2\pi b(r')r'^2[T_{rr}(r')+S(r')]\,\dd r'. \tag{IEF35}
\]
This outward solve infers the Arnowitt--Deser--Misner (ADM) mass.  After that,
or whenever a black hole interior is represented, the correct formulation is
the fixed-mass inward solve
\[
 Y(r)=M_{\rm ADM}
 -\int_r^\infty2\pi b(r')r'^2[T_{rr}(r')+S(r')]\,\dd r'. \tag{IEF36}
\]
On a finite vacuum outer boundary \(R_E\),
\[
 b(R_E)=1,\qquad Y(R_E)=M_{\rm ADM}.                      \tag{IEF37}
\]
Production fresh EF runs do not use the regular-center inference path.  They
authenticate a solved null initial-data artifact, import its
\(M_{\rm null}\), and set \(M_{\rm ADM}=M_{\rm null}\) from the first EF field
solve onward.  Repeating the outward regular-center solve after excision would
erase the hole mass and is forbidden.

\subsubsection{Vlasov equation and particle characteristics}
\label{sec:ef-vlasov-derivation}

The distribution function
\[
                 f=f(v,r,P,L^2)
\]
is a scalar on the future unit-mass shell.  Its arguments are advanced time
\(v\), areal radius \(r\), the covariant canonical radial momentum
\(P:=p_r>0\), and the conserved squared angular momentum magnitude \(L^2\).
VENUS stores \(\eta:=L^2/2\) and uses
\[
 \chi(r,\eta):=1+\frac{2\eta}{r^2}.                       \tag{IEF9}
\]
The symbol \(p_a\) denotes the particle four-momentum covector; unit particle
rest mass is used, while a marker's invariant rest-number weight represents
many physical particles.

The mass-shell condition \(g^{ab}p_a p_b=-1\) is
\[
 -1=\frac{2}{b}p_v P+aP^2+\frac{2\eta}{r^2}.             \tag{IEF11}
\]
Because \(g^{vr}=1/b\),
\[
 p^v=\frac{P}{b}>0,                                      \tag{IEF10}
\]
so \(P>0\) selects the future branch; it does not mean outward motion.  Solving
the mass shell for the momentum conjugate to \(v\) gives the advanced-time
Hamiltonian
\[
 p_v=-\frac b2\left(aP+\frac{\chi}{P}\right),\qquad
 {\cal H}_{\rm EF}:=-p_v
 =\frac b2\left(aP+\frac{\chi}{P}\right).                \tag{IEF12}
\]
The contravariant radial momentum is
\[
 p^r=\frac12\left(aP-\frac{\chi}{P}\right).               \tag{IEF13}
\]

Collisionless transport means that \(f\) is constant along geodesic
characteristics:
\[
 \partial_v f+\dot r\,\partial_r f+\dot P\,\partial_P f=0,
 \qquad \dot\eta=0,                                      \tag{IEF14}
\]
where a dot now means \(\dd/\dd v\).  Hamilton's equations give every
characteristic:
\[
 \boxed{\dot r=\frac{\partial{\cal H}_{\rm EF}}{\partial P}
 =\frac b2\left(a-\frac{\chi}{P^2}\right),}              \tag{IEF15}
\]
\[
 \boxed{\dot P=-\frac{\partial{\cal H}_{\rm EF}}{\partial r}
 =-\frac P2(ab_r+ba_r)-\frac{b_r\chi}{2P}
   +\frac{2b\eta}{P r^3}.}                               \tag{IEF16}
\]
In the last derivative \(P\) and \(\eta\) are held fixed.  The first term
differentiates the metric factor \(ab\) multiplying the radial kinetic
momentum.  The second differentiates the factor \(b\) multiplying rest and
angular energy.  The last follows from
\(\partial_r\chi=-4\eta/r^3\) and is the spherical centrifugal term.  These
are coordinate forms of geodesic motion, not separate forces added to the
Einstein equations.

The sign of physical radial motion is
\[
 \dot r<0\Longleftrightarrow aP^2<\chi,\qquad
 \dot r=0\Longleftrightarrow aP^2=\chi,\qquad
 \dot r>0\Longleftrightarrow aP^2>\chi.                  \tag{IEF13a}
\]
At an apparent horizon,
\[
 \left.\dot r\right|_{a=0}=-\frac{b\chi}{2P^2}<0.         \tag{IEF17}
\]
A future massive particle therefore crosses inward at finite advanced time.
No \(1/a\) occurs in (IEF15)--(IEF16).  A step that loses \(P>0\) is invalid
and is rejected rather than repaired with a floor.

\paragraph{Term-by-term correspondence with the SC characteristics.}
Let \(w\) be SC orthonormal radial momentum and define the local energy per
unit rest mass
\[
 \epsilon:=\sqrt{1+w^2+\frac{2\eta}{r^2}}.                \tag{IEF41}
\]
The SC characteristic equations are
\[
 \frac{\dd r}{\dd\tau}=e^{\mu-\lambda}\frac{w}{\epsilon},
 \qquad
 \frac{\dd w}{\dd\tau}
 =-\lambda_\tau w-e^{\mu-\lambda}\mu_r\epsilon
  +e^{\mu-\lambda}\frac{L^2}{r^3\epsilon}.               \tag{IEF-SC-V}
\]
Their three momentum terms describe the changing radial metric, the lapse
gradient, and the centrifugal barrier.

A momentum is a covector, so equality of the canonical one-form at the same
event,
\[
 p_\tau\dd\tau+p_r^{({\rm SC})}\dd r=p_v\dd v+P\dd r,
\]
with \(p_\tau=-e^\mu\epsilon\) and
\(p_r^{({\rm SC})}=e^\lambda w\), gives
\[
 \boxed{P=e^\lambda(\epsilon+w)},\qquad
 \boxed{-p_v=F e^\mu\epsilon}.                            \tag{IEF42}
\]
The subtraction-free expression
\[
 P=\frac{1+2\eta/r^2}{\sqrt a\,(\epsilon-w)}              \tag{IEF43}
\]
is used for strongly infalling data.  Along a particle,
\[
 \frac{\dd v}{\dd\tau}
 =A\left(1+\frac{w}{\epsilon}\right)>0.                  \tag{IEF42a}
\]
Dividing the SC radial equation by this positive factor gives (IEF15)
directly.

For the momentum equation set
\(x:=\epsilon+w=\sqrt a\,P\) and \(c=e^{\mu-\lambda}\).
Differentiating \(P=e^\lambda x\) along the SC characteristic and dividing by
\(\dd v/\dd\tau\) yields
\[
 \frac{\dot P}{P}
 =F\left[
   \lambda_\tau\frac{\epsilon-w}{x}
  +c\lambda_r\frac{w}{x}
  -c\mu_r\frac{\epsilon}{x}
  +c\frac{L^2}{r^3x^2}\right].                           \tag{IEF-SC-V1}
\]
Substitution of the metric identities in
(IEF38d), (IEF-SC-E1), and (IEF-SC-E2) reduces this expression exactly to
(IEF16).  The four displayed pieces are respectively the SC time-dependent
radial metric term, spatial radial-metric term, lapse force, and angular
barrier.  Their regrouping into \(a_r\), \(b_r\), and \(\chi_r\) explains
every EF term.

\paragraph{Stress energy carried by the distribution.}
Using \(\eta=L^2/2\), the independent covariant components are
\begin{align}
 T_{vv}
 &=\frac{2\pi}{r^2}\int_0^\infty\!\int_0^\infty
   \frac{p_v^2}{P}f\,\dd\eta\,\dd P,                  \tag{IEF18}\\
 T_{vr}
 &=\frac{2\pi}{r^2}\int_0^\infty\!\int_0^\infty
   p_v f\,\dd\eta\,\dd P,                            \tag{IEF19}\\
 T_{rr}
 &=\frac{2\pi}{r^2}\int_0^\infty\!\int_0^\infty
   P f\,\dd\eta\,\dd P,                              \tag{IEF20}\\
 T_{\theta\theta}
 &=\frac{2\pi}{r^2}\int_0^\infty\!\int_0^\infty
   \frac{\eta}P f\,\dd\eta\,\dd P,                \tag{IEF21}\\
 T_{\phi\phi}&=\sin^2\theta\,T_{\theta\theta}.           \tag{IEF22}
\end{align}
Using (IEF12), the constraint source introduced in (IEF23a) is
\[
 S=\frac{2\pi}{r^2}\int_0^\infty\!\int_0^\infty
   \frac{1+2\eta/r^2}{P}f\,\dd\eta\,\dd P.         \tag{IEF23}
\]
Thus \(T_{rr}\) weights radial canonical momentum, while \(S\) weights rest
and angular energy divided by that momentum.  VENUS deposits precisely these
two combinations for the radial constraint solve; it reconstructs the other
components for profiles and coordinate-parity checks.

\subsubsection{What the production code solves}
\label{sec:ef-what-code-solves}

A production EF step does not freely evolve all four Einstein components.
At each advanced time it deposits the active marker ensemble, solves the two
radial constraints (IEF28)--(IEF29), and advances the Vlasov characteristics
(IEF15)--(IEF16).  The remaining Einstein equations are consistency
relations.  WPIC means weighted particle in cell: each numerical marker has
an invariant rest-number weight, and its cloud deposits a fraction of that
weight into neighboring radial cells.

Constructing an ingoing null initial slice requires the SC solution at the
different times where particle worldlines cross that slice.  The version-2
donor, abbreviated donor-v2, stores this dense SC history together with the
center and face data needed to reconstruct its metric.  The solved null
artifact is the binary initial-data file obtained from those crossings.  It
contains the mapped markers, metric and stress data, boundary mass, accuracy
record, and hashes that identify its inputs.
The ledger is the ordered list of every accepted donor step and its exact
timestep, used for deterministic replay.  A manifest is a JavaScript Object
Notation (JSON) metadata file that names the artifacts and their 256-bit
Secure Hash Algorithm (SHA-256) digests.  Provenance means this recorded chain
from physical initial condition and source version to each derived object.
A check is called fail-closed when complete, consistent evidence is required
to pass.  Missing or inconsistent data then produces a stop or an
inconclusive result.

The public data objects used below have distinct roles.  A \code{Marker}
stores invariant rest weight, \(r,P,\eta\), persistent and source identifiers,
and an active, excised, or escaped state.  \code{DepositedSources} stores cell
and shell source arrays.  \code{Metric} stores center and face values of
\(a,b,m,a_r,b_r\), while \code{Field} pairs that metric with its deposited
sources.  \code{ConstraintOptions} selects the radial mass boundary and
excision controls; \code{EvolutionOptions} adds the pusher, particle shape,
thread count, and timestep controls.  A \code{StepResult} reports the accepted
step, endpoint field, fixed mass, and excision outcome.  A map
\code{EikonalSurface} stores the solved null-surface event times, Jacobians,
and refinement evidence.

\begin{figure}[ht]
\centering
\textbf{Physics flow: ingoing-EF Einstein--Vlasov WPIC}
\par\smallskip
\begin{tikzpicture}[node distance=0.22cm]
\node[physicsflow] (ef1) {Authenticate the solved null initial-data artifact, its donor-v2 history, and the fixed mass \(M_{\rm null}\)};
\node[physicsflow,below=of ef1] (ef2) {Deposit active markers into the null-slice sources \(T_{rr}\) and \(S\)};
\node[physicsflow,below=of ef2] (ef3) {Integrate \(b\) and \(mb\) inward from the vacuum outer boundary; form \(a\), metric derivatives, and the apparent-horizon diagnostic};
\node[physicsflow,below=of ef3] (ef4) {Choose a causality and momentum limited advanced-time step};
\node[physicsflow,below=of ef4] (ef5) {Push a copy of the markers through half a step in the old field};
\node[physicsflow,below=of ef5] (ef6) {Redeposit that copy, solve the midpoint field, and push the original markers through the full step};
\node[physicsflow,below=of ef6] (ef7) {Solve the endpoint field, apply only causally available excision, and recompute the fixed-mass endpoint field};
\node[physicsflow,below=of ef7] (ef8) {Audit every accepted state; write diagnostics and an atomic restart; then repeat};
\draw[physicsarrow] (ef1) -- (ef2);
\draw[physicsarrow] (ef2) -- (ef3);
\draw[physicsarrow] (ef3) -- (ef4);
\draw[physicsarrow] (ef4) -- (ef5);
\draw[physicsarrow] (ef5) -- (ef6);
\draw[physicsarrow] (ef6) -- (ef7);
\draw[physicsarrow] (ef7) -- (ef8);
\end{tikzpicture}
\caption{Code-logic flow for one production ingoing-EF Einstein--Vlasov
WPIC update.  The fixed null-slice mass is retained when interior marker
sources are excised.}
\label{fig:ef-code-logic-flowchart}
\end{figure}

\paragraph{Initial data and first field solve.}
A fresh production EF run is allowed only in the high-compactness scenario,
on a linear radial grid with the configured origin-parity mesh convention.
It must receive both the authenticated coordinate-neutral physical initial
condition and a solved EF null artifact.  The artifact records the
worldline/null-surface intersections, EF marker momenta, reconstructed
center/face fields, map errors, donor ledger, hashes, and \(M_{\rm null}\).
Supplying only a spacelike SC marker file is a hard error.  The driver verifies
the physical-initial-condition hash, base and realized fingerprints, sampling
realization, marker count and grid, map-manifest hash, artifact hashes, and
initial metric/stress pullbacks before evolution.  Its first and all later
constraint solves use the inward boundary condition
\(Y(R_E)=M_{\rm null}\).  The regular-center mode (IEF35) remains available
only as a reusable library mode and unit-test target.

A restart is an atomic checkpoint containing the state needed to continue the
same evolution.  A resume reads the EF markers, advanced time, step number,
excision mode, and fixed boundary mass, then reconstructs the metric from
those markers. When physical or null initial data are supplied, it compares
each marker's invariant rest weight, angular momentum, identifier, and source
identifier with that authenticated realization before opening output files.
An accepted-step audit sidecar is trusted only when its exact
step number and advanced time match the restart.  A missing, stale, or
old-schema sidecar does not invent evidence: evolution may continue, but
black hole certification remains disabled until a complete new audit exists.
Here a sidecar is a separate atomic evidence file tied to the checkpoint; a
run summary is the machine-readable record of a completed run, and a
certificate is a separately authenticated record that promotes a scientific
claim only after all named checks pass.
Atomic means that the completed file replaces the old file as one operation,
so an interrupted write is not mistaken for a valid checkpoint.

\paragraph{Boundary, horizon, and excision conditions.}
The map worldtube \(R_B\) and the EF evolution boundary \(R_E\) are distinct.
The extraction requires \(R_B\le R_E\); the initial null surface between them
may use the audited analytic vacuum extension.  Evolution uses the full
domain through \(R_E\).  At the outer boundary,
\(b(R_E)=1\) and \(Y(R_E)=M_{\rm null}\).  This assumes the boundary stays in
vacuum; fixed total mass cannot reconstruct the finite-radius field of
dynamically important matter that has left the grid.
An escaped marker is an active marker that leaves the evolution domain; its
record and rest weight remain in accounting, but it no longer deposits on the
finite grid.

At the inner side, the origin-parity option creates a true \(r=0\) field face
but does not reflect EF characteristics.  Before an excision surface exists,
an active marker must remain at positive radius inside the field domain.
An attempted inner crossing rejects the trial step; the interior-evolution driver
retries with a smaller timestep as described below.
The characteristic equations are singular at \(r=0\) for \(\eta>0\).

The metric solver searches face intervals for the outermost
negative-to-positive zero of \(a\) and interpolates \(r^2a\) to obtain the
apparent-horizon radius \(r_H\), with \(m_H=r_H/2\).  Excision is enabled only
after a cell center satisfies \(a<a_{\rm ex}<0\); the particle boundary is
placed a configured number of cell centers farther inward and at most a fraction
\(F=0.1\) of the apparent-horizon radius.  The selected center, every center in the outward buffer, and the intervening
faces must satisfy \(a<a_{\rm ex}\), with finite positive \(b\).  A trapped outer
shell surrounding an untrapped central cavity therefore supplies a boundary
only when a qualifying trapped center lies within this radius cap.  Markers at or
inside it change from active to excised, keep their records and invariant
weights, stop depositing, and are no longer pushed.  The endpoint field is
then recomputed with fixed \(M_{\rm null}\).  No inner boundary data are
supplied inside this trapped, causally disconnected region.
An excision radius already present in the old field is enforced during every
frozen-metric push.  An endpoint-supported candidate boundary
is selected from the recomputed trapped region; it is applied before the step is
accepted and that endpoint field is recomputed. The retained boundary is
checked in this recomputed metric, including its buffer and radius cap.

The option \code{--ef-evolve-interior} retains all interior markers. Horizon
detection continues, while every metric stage sets the applied excision
radius to zero. The markers continue to move and deposit their stress.
Section~\ref{sec:ef-interior-evolution} gives the timestep, restart, and resolution
conditions for this mode.

\paragraph{Time integration, rejection, and accepted-state audit.}
The production driver uses two equal characteristic substeps and
\(C_{\rm CFL}=0.20\), where the Courant--Friedrichs--Lewy (CFL) number limits
radial travel and momentum change per step.  Internally the RK4 pusher evolves
\(u=\log P\), which preserves the mathematical \(P>0\) branch without imposing
a momentum floor.  A nonfinite characteristic stage, invalid logarithmic
state, or retryable right-hand-side (RHS) or RK4 failure rejects the whole
self-consistent step;
the driver halves \(\Delta v\) and retries, at most 24 times. In interior-evolution
mode, an attempted inner crossing or a nonfinite trial constraint solve also
allows this retry. The last accepted state is saved when the available
retries or floating-point time resolution are exhausted in interior-evolution mode.

Before output cadence is considered, every accepted state is audited for
finite markers, sources, center and face metric arrays, \(P>0\), \(b>0\),
fixed-\(M_{\rm null}\) agreement, causal excision, and positive excised weight
when excision is claimed.  Saved history then records marker-weight channels,
ADM mass, compactness, minimum \(a\) and \(b\), horizon and excision radii,
Lagrangian radii, virial and Jeans diagnostics.  Full stress profiles support
independent coordinate pullbacks.  General independent residuals of
(IEF28)--(IEF31) and a generic evolved-exterior vacuum norm are recommended
validation diagnostics, but are not currently emitted as ordinary
production-history columns; the implemented exact-vacuum residual belongs to
the authenticated null-map extension audit.
\subsubsection{Detailed coordinate map, donor-v2 history, and fail-closed provenance}
\label{sec:ef-detailed-map}

Section~\ref{sec:ef-eikonal-construction} established the geometry and the
forward dynamical map.  This section records the component identities and
numerical evidence used by the implementation.  It is deliberately placed
after the field and particle equations because the map initializes those
equations and identifies the physical events used to compare them.

At one mapped event the covariant and contravariant radial metric blocks are
\begin{align}
 g^{\rm EF}_{vv}&=-ab^2=-F^2e^{2\mu},&
 g^{\rm EF}_{vr}&=b=F e^{\lambda+\mu},&
 g^{\rm EF}_{rr}&=0,                                    \tag{IEF38d-1}\\
 g_{\rm EF}^{vv}&=0,&
 g_{\rm EF}^{vr}&=b^{-1}=A e^{-\lambda-\mu},&
 g_{\rm EF}^{rr}&=a=e^{-2\lambda}.                       \tag{IEF38d-2}
\end{align}
These are component transformations, not an identification of the raw arrays
\((\lambda,\mu)\) and \((a,b)\).  The retarded root reverses the mixed term
and has \(b_{\rm out}<0\); the EF evolution namespace implements only the
future advanced branch \(b>0\).

For a static exterior, both directions of the transformation are
\begin{align}
 \dd v&=\dd\tau+e^{\lambda-\mu}\dd r
       =\dd\tau+\frac{\dd r}{ab},
 &v(\tau,r)&=\tau+\int^r\frac{\dd r'}{a(r')b(r')},       \tag{IEF40a}\\
 \dd\tau&=\dd v-\frac{\dd r}{ab},
 &\tau(v,r)&=v-\int^r\frac{\dd r'}{a(r')b(r')}.          \tag{IEF40b}
\end{align}
The corresponding metric functions satisfy
\[
 e^{2\lambda}=\frac1a,\qquad e^{2\mu}=ab^2,\qquad
 \lambda=-\frac12\log a,\qquad
 \mu=\log b+\frac12\log a.                               \tag{IEF40c}
\]
For a dynamical inverse map, the integrating factor in (IEF40d) must satisfy
\[
 \partial_r F=-\partial_v\left(\frac{F}{ab}\right),       \tag{IEF40e}
\]
and substitution gives
\[
 \dd s^2=-\frac{ab^2}{F^2}\dd\tau^2
          +\frac1a\dd r^2+r^2\dd\Omega^2.                \tag{IEF40f}
\]
Setting \(F=1\) in a time-dependent field would impose the extra restriction
\(\partial_v[(ab)^{-1}]=0\); it is not a freely available gauge choice.
Neither \(F\) nor the event correspondence can be reconstructed from one
metric snapshot.

All SC-to-EF transformations and their inverse checks are restricted to the
connected overlap with \(a>0\).  As \(a\to0^+\), \(e^{2\lambda}=1/a\) and
\(1/(ab)\) diverge, so SC data and the public momentum-map helpers stop.
Finite positive \(b\) leaves the EF metric and EF characteristics regular;
the restriction belongs to the comparison chart, not the EF evolution.

\paragraph{Momentum components carried by the mapped artifact.}
If \(w\) is the existing VENUS orthonormal radial momentum and
\[
 \epsilon=\sqrt{1+w^2+\frac{2\eta}{r^2}},
\]
then the polar-coordinate covector components and inverse-coordinate
derivatives at the same event are
\[
 p_\tau=-e^\mu\epsilon,\qquad
 p_r^{({\rm pol})}=e^\lambda w,\qquad
 \left.\partial_v\tau\right|_r=F,\qquad
 \left.\partial_r\tau\right|_v=-\frac{F}{ab}
 =-e^{\lambda-\mu}.                                    \tag{IEF41a}
\]
A momentum is a one-form,
\[
 p=p_\tau\dd\tau+p_r^{({\rm pol})}\dd r
   =p_v\dd v+P\dd r.                                  \tag{IEF41b}
\]
Its components transform with derivatives of the old coordinates.  Using
(IEF40d), and equivalently \(ab=F e^{\mu-\lambda}\), gives
\[
 \boxed{
 P=p_r^{({\rm pol})}+p_\tau
       \left.\partial_r\tau\right|_v
  =e^\lambda(w+\epsilon)
  =\frac{w+\epsilon}{\sqrt a}},\qquad
 p_v=p_\tau\left.\partial_v\tau\right|_r
    =-F e^\mu\epsilon.
\]
The integrating factor cancels from \(P\).  Hence (IEF42) remains an exact
pointwise identity for a dynamical coordinate map when its SC and EF
arguments describe the same event.  The helper needs \(a\), but not \(F\);
it does not determine which events on two different slices correspond.
For strongly infalling data the subtraction-free evaluation is
\[
 P=\frac{1+2\eta/r^2}{\sqrt a\,(\epsilon-w)}.
\]
The inverse map, defined only for \(a>0\), is
\[
 x:=\sqrt a\,P,\qquad
 w=\frac12\left(x-\frac{\chi}{x}\right),\qquad
 \epsilon=\frac12\left(x+\frac{\chi}{x}\right).       \tag{IEF44}
\]
The same covector transformation gives
\[
 \frac b2\left(aP+\frac\chi P\right)
 =-p_v=F e^\mu\epsilon.                               \tag{IEF45}
\]
Thus the right-hand side is \(e^\mu\epsilon\) in the static normalization
\(F=1\).  The factor \(F\) is essential in a general dynamical
polar-to-advanced coordinate transformation; it is already encoded in the EF
metric function \(b=F e^{\lambda+\mu}\).

For reference, the retarded/outgoing coordinate and its future-directed
signed momentum are
\[
 u_{\rm ret}=\tau-\int^r e^{\lambda-\mu}\dd r',\qquad
 b_{\rm out}=-e^{\lambda+\mu},\qquad
 (p_r)_{\rm out}=e^\lambda(w-\epsilon)<0.              \tag{IEF45a}
\]
East's physical-data recipe first evolves an outward/dispersing solution and
then uses \(t\mapsto-t\) to obtain its incoming time reverse.  With
\(w_{\rm in}=-w_{\rm out}\), it acts as
\[
 b_{\rm in}=-b_{\rm out}>0,\qquad
 (p_r)_{\rm in}=-(p_r)_{\rm out}>0,\qquad
 r_{\rm in}=r_{\rm out},\qquad
 \eta_{\rm in}=\eta_{\rm out}.                        \tag{IEF45b}
\]
and therefore
\[
 -(p_r)_{\rm out}
 =e^\lambda(\epsilon-w_{\rm out})
 =e^\lambda(\epsilon+w_{\rm in})=(p_r)_{\rm in}.       \tag{IEF45c}
\]
Time reversal here means more than replacing the printed time coordinate by
its negative: every radial velocity or momentum is reversed as well, while
positions, angular momentum magnitudes, and invariant weights are retained.
Changing \(t\) without changing \(w\) would not construct the incoming
solution.
VENUS does not use that operation to relabel an outgoing production run as its
physical collapse.  Its physical neutral shell receives the controlled inward
kick (IEF-S26), so the physical Schwarzschild evolution is already ingoing and
is never time-reversed.  The null-slice extractor nevertheless needs the part
of that physical solution at negative Schwarzschild times because a single
ingoing null slice intersects different radii at different \(\tau\).  It
obtains that history without backward numerical integration by making an
\emph{independent} donor copy and using East's construction exactly once:
\[
 s=-\tau,\qquad w_{\rm reverse}(s=0)=-w_{\rm physical}(\tau=0).
                                                               \tag{IEF45c-1}
\]
The copy is evolved forward in \(s\).  At a marker--null-surface crossing the
physical momentum is restored as
\(w_{\rm physical}=-w_{\rm reverse}\), after which (IEF42) gives the
future-directed \(P>0\).  Thus ``using \(t\mapsto-t\)'' refers only to this
auxiliary donor/history construction.  The incoming Schwarzschild run is
unchanged, and the EF characteristic pusher never reverses its time or its
momentum.

The public conversion routines implement (IEF42)--(IEF44), including the
stable infalling branch:
\codesource{SC--EF marker-momentum conversion}{\path{cpu/include/ev/weighted_einstein_vlasov_pic_eddington_finkelstein.hpp}}{Public same-event momentum-map helpers and EF marker construction}
\begin{lstlisting}[language=C++,basicstyle=\ttfamily\scriptsize]
namespace ef = ev::eddington_finkelstein;

const double p_r = ef::canonical_p_r_from_schwarzschild_w(
    marker.w, marker.eta, marker.r, metric_gap);
ef::Marker ef_marker = ef::from_schwarzschild_marker(
    marker, metric_gap);

const double w = ef::schwarzschild_w_from_canonical_p_r(
    ef_marker.p_r, ef_marker.eta, ef_marker.r, metric_gap);

const double p_r_out =
    ef::outgoing_canonical_p_r_from_schwarzschild_w(
        outgoing_w, eta, radius, metric_gap);
const double p_r_in =
    ef::time_reverse_outgoing_canonical_p_r(p_r_out);
\end{lstlisting}
The marker rest weight and \(\eta\), its persistent identifier (ID), and its
source identifier are unchanged by this conversion.

The production extractor makes the separation between physical and reverse
states explicit (irrelevant checks are elided here):
\codesource{Solved-null initial-data extraction}{\path{cpu/src/drivers/einstein_vlasov_wpic.cpp}}{Production reverse-donor copy, eikonal solve, deterministic replay, and marker crossings}
\begin{lstlisting}[language=C++,basicstyle=\ttfamily\scriptsize]
std::vector<ev::WeightedPICMarker> reverse_initial = physical_markers;
for (auto& marker : reverse_initial) {
    marker.w = -marker.w;                 // s=-tau donor only
}
std::vector<ev::WeightedPICMarker> donor_markers = reverse_initial;
// Store every accepted donor field and the exact accepted-dt ledger.
map::EikonalSurface surface = map::solve_ingress_eikonal_surface(
    history, grid, reverse_time, 8);
// Replay the ledger bit-for-bit and require one crossing for every ID.
crossings.push_back(map::interpolate_marker_crossing(
    before[q], replay[q], replay_time, next_time,
    surface, history, grid, step + 1));
\end{lstlisting}
The integer passed to the surface solver is \(n\), the number of RK4
substeps used for the coarse solve on each radial cell.  Starting from the
same already accepted state \(Y_{i,\mathrm{in}}\) at the outer end of cell
\(i\), the extractor independently computes an \(n\)-substep result
\(Y_i^{(n)}\) and a \(2n\)-substep result \(Y_i^{(2n)}\), where

\[
        Y=(s_\sigma,\log A).
\]
It accepts \(Y_i^{(2n)}\) as the incoming state for the next cell; it does not
advance with the coarse result and does not Richardson-extrapolate the state.
For each component it estimates the local error of that accepted, refined
fourth-order result by
\begin{align}
 \epsilon_{s,i}
 &=\frac{\left|s_{\sigma,i}^{(2n)}-s_{\sigma,i}^{(n)}\right|}
 {\left(2^4-1\right)
  \max\!\left(1,
       \left|s_{\sigma,i,\mathrm{in}}\right|,
       \left|s_{\sigma,i}^{(2n)}\right|\right)},                 \tag{IEF45c-2}\\
 \epsilon_{\log A,i}
 &=\frac{\left|\left(\log A\right)_i^{(2n)}
                    -\left(\log A\right)_i^{(n)}\right|}
 {\left(2^4-1\right)
  \max\!\left(1,
       \left|\left(\log A\right)_{i,\mathrm{in}}\right|,
       \left|\left(\log A\right)_i^{(2n)}\right|\right)} .       \tag{IEF45c-3}
\end{align}
The unit term gives a dimensionless error that behaves as an absolute error
for small values and a relative error for large ones.  It uses one code-time
unit for \(s_\sigma\).  The factor \(2^4-1=15\) is the RK4
step-doubling Richardson divisor.  The surface records the two maxima
separately as
\texttt{maximum\_reverse\_time\_richardson\_error} and
\texttt{maximum\_log\_A\_richardson\_error}.  The authenticated manifest
serializes them unambiguously as
\path{maximum_eikonal_reverse_time_richardson_error} and
\path{maximum_eikonal_log_A_richardson_error}, together with the exact
normalization string in (IEF45c-2)--(IEF45c-3).  Its backward-compatible
aggregate is exactly
\[
 \texttt{maximum\_eikonal\_ode\_residual}
 =\max_i\left\{\epsilon_{s,i},\epsilon_{\log A,i}\right\}.       \tag{IEF45c-4}
\]
Despite its retained manifest name, (IEF45c-4) is an integration-error
estimate, not the finite-difference quantity

\[
 \left|\frac{\Delta s_\sigma}{\Delta r}
              -h(s_{\mathrm{mid}},r_{\mathrm{mid}})\right| .     \tag{IEF45c-5}
\]
The latter includes the midpoint quadrature defect produced by the curvature
of the physical null ray.  It can therefore remain larger than the RK4
integration error even when the ray is accurately integrated.  It is not used as an acceptance residual.
The null-map preflight applies the same, unchanged \(2\times10^{-5}\)
ceiling to each authenticated extraction component and to their aggregate.
This integration-error estimate separates RK4 error from the physical
curvature of the ray while preserving the acceptance tolerance.  The later same-event history comparator solves a
different initial-value problem for each sampled event.  It deliberately
uses the still more conservative raw \(n\)-versus-\(2n\) change described
below, rather than assuming that a ray crossing time-interpolation knots is
already in the smooth fourth-order Richardson regime.

The binary null slice and its manifest, crossing table, mapped field slice,
accepted-step ledger, and dense donor history are all recorded with the
256-bit Secure Hash Algorithm (SHA-256).  A
fresh EF run refuses to start without this solved-eikonal artifact and verifies
the base initial condition (IC), realized Schwarzschild IC, sampling
realization, grid, angular
controls, marker identities, and all nested artifact hashes.

\paragraph{Separated map worldtube and evolution boundary.}
The outer radius is part of the characteristic initial-data problem, not a
harmless output choice.  VENUS therefore distinguishes the exact mesh-face
worldtube \(R_B\), set by \code{--ef-map-worldtube-radius}, where the numerical
eikonal solve is normalized, from the full EF evolution boundary \(R_E\), set
by \code{--high-r-max}.  They obey
\[
 0<R_B\leq R_E,\qquad R_B>2M_{\rm null}.                \tag{IEF45c-6}
\]
The legacy case is \(R_B=R_E\).  When \(R_B<R_E\), every matter worldline has
already crossed the selected null surface strictly inside \(R_B\), so the
remainder of that \emph{initial null slice} is known analytically.  Define
\[
 \Delta r_\star(r;R_B,M)=
 \begin{cases}
 r-R_B, &M=0,\\[2pt]
 (r-R_B)+2M\log\!\left(\dfrac{r-2M}{R_B-2M}\right),&M>0,
 \end{cases}                                           \tag{IEF45c-7a}
\]
which is \(r_\star(r)-r_\star(R_B)\) without a logarithm of a dimensionful
quantity.  If \(s=-\tau\), and \(s_B\) is the reverse-donor time at the
worldtube, the exact extension for \(R_B<r\leq R_E\) is
\begin{align}
 s(r)&=s_B+\Delta r_\star(r;R_B,M_{\rm null}),&
 \tau(r)&=-s_B-\Delta r_\star(r;R_B,M_{\rm null}),
                                                               \tag{IEF45c-7b}\\
 m(r)&=M_{\rm null},&
 a(r)&=1-\frac{2M_{\rm null}}r,                       \tag{IEF45c-7c}\\
 A(r)&=F(r)=b(r)=1,&
 T_{ab}(r)&=S(r)=0.                                    \tag{IEF45c-7d}
\end{align}
In particular, \(F=1\), not \(1/a\).  The radial derivative of the inverse
time map at fixed \(v\) is instead
\[
 \left.\frac{\partial\tau}{\partial r}\right|_v
 =-\frac{F}{ab}=-\frac1a,                              \tag{IEF45c-7e}
\]
exactly as the vacuum tortoise relation requires.  The mass in this extension
is \(M_{\rm null}\), the ADM mass on the extracted null slice; using
\(M_{\rm family}\) or the independently deposited \(M_{\rm dep}\) would
introduce a boundary inconsistency.

The eikonal construction labels the extracted surface by \(v=-T_{\rm donor}\)
at the normalized worldtube.  Runtime output shifts only the additive time
origin:
\[
 \boxed{\ \widehat v:=v+T_{\rm donor}\ },
 \qquad \widehat v=0\quad\hbox{on the imported initial null slice}.          \tag{IEF45c-7f}
\]
The manifest names the same quantity
\(\code{v_bar=v+ef_v_offset}\), with
\(\code{ef_v_offset}=T_{\rm donor}\).  No metric or momentum component is
changed by this constant relabelling.

This split is accepted only after all of the following fail-closed gates:
\begin{enumerate}
\item \(R_B\) is an exact non-inner edge of the full \(R_E\) mesh, so the
      numerically mapped cells form one authenticated grid prefix;
\item the last four donor cells below \(R_B\) have vanishing matter sources,
      the discrete vacuum value \(e^{\lambda+\mu}=1\) satisfies its recorded
      \(O(\Delta r_B^2)\) tolerance, and every donor worldline remains at
      least four local cells inside \(R_B\);
\item every persistent marker ID has one and only one surface crossing, each
      crossing has four-cell clearance from \(R_B\), and no marker reaches
      the finite evolution boundary first;
\item the complete numerical surface lies inside the dense donor history,
      \(A>0\), \(F>0\), \(a>0\) throughout the overlap, and the two separate
      RK4 step-doubling errors in (IEF45c-2)--(IEF45c-4) pass;
\item the analytic part has zero source, constant \(m=M_{\rm null}\),
      \(A=F=b=1\), positive \(a\), and vanishing field and inverse-map
      residuals at their recorded tolerance.
\end{enumerate}
The version-2 binary and JavaScript Object Notation (JSON) manifest record
\(R_B\), \(R_E\), the analytic
extension start index, separate prefix and full-grid fingerprints, both
clearances, vacuum-source and \(b\)-normalization audits, and all analytic
extension residuals.  These metadata and the slice, crossing table,
accepted-step ledger, donor history, physical IC, and nested hashes are
authenticated again on EF import and resume.

The analytic continuation in (IEF45c-7a)--(IEF45c-7e) constructs only the
vacuum part of the \emph{initial} characteristic surface.  It does not freeze
the region \(R_B<r<R_E\) during evolution: after import the EF constraints and
characteristics are evolved on the complete \(R_E\) grid.  Nor do values of
\(R_B,R_E,T_{\rm donor}\), marker count, or grid resolution transfer from an
older Gaussian-shell experiment.  The bound static-shell family below must
derive all of them from its own local crossing/vacuum preflight before any
long job is submitted.

\paragraph{Slice caveat.}
Equations (IEF39)--(IEF44) are an algebraic/radial map only for a static
solution.  Dynamically, (IEF38a)--(IEF40e) must first determine which events
correspond.  The pointwise momentum identity (IEF42) remains exact at those
events, but the pointwise conversion helpers in the EF kernel neither solve
the eikonal PDE nor resample a hypersurface; the separate authenticated
coordinate-map pipeline described below performs those operations.  A
\(\tau={\rm const}\) surface is spacelike whereas a
\(v={\rm const}\) surface is null, and \(v=0\) generally samples different SC
times at different radii.  A rigorous cross-coordinate validation should
therefore use a transformed static state, an analytic worldline, or
the same physical events in a dynamical spacetime.  Merely copying the same radial arrays
into both modules does not create identical physical initial hypersurfaces.

\paragraph{What equality of the two coordinate descriptions means.}
This point is essential for interpreting any SC--EF comparison.  The
distribution function is a scalar on the future mass shell, not an array
whose entries may be compared at the same pair of numerical time and radial
indices.  At one and the same event, with the momenta related by
(IEF42)--(IEF44), scalar invariance is the statement
\[
 f_{\rm EF}\bigl(v(\tau,r),r,P,\eta\bigr)
 =f_{\rm SC}\bigl(\tau,r,w,\eta\bigr).             \tag{IEF45d}
\]
There is no multiplicative Jacobian in (IEF45d).  A histogram density per
coordinate momentum interval does carry one: if
\(n_w\,\dd w=n_P\,\dd P\), then
\[
 n_P=n_w\left|\frac{\partial w}{\partial P}\right|,
 \qquad
 \frac{\partial w}{\partial P}
 =\frac{\sqrt a}{2}\left(1+\frac{\chi}{aP^2}\right). \tag{IEF45e}
\]
WPIC marker rest weights are invariant particle-number weights and therefore
remain unchanged; (IEF45e) is relevant only when the empirical measure is
converted into a density in a chosen momentum coordinate.
The forward Jacobian and the reduced invariant mass-shell measure make this
distinction explicit:
\[
 \frac{\partial P}{\partial w}=\frac P\epsilon>0,
 \qquad
 \boxed{\;\frac{\dd P}{P}=\frac{\dd w}{\epsilon}\;},
 \qquad
 \dd\Pi_{\rm red}
 =\frac{2\pi}{r^2}\frac{\dd w\,\dd\eta}{\epsilon}
 =\frac{2\pi}{r^2}\frac{\dd P\,\dd\eta}{P}.      \tag{IEF45e-1}
\]
Thus substituting the scalar relation (IEF45d) and (IEF45e-1) into the
stress integrals reproduces the same tensor without an ad hoc rescaling of
\(f\).

A finite WPIC realization does not provide a pointwise value of the smooth
continuum function \(f\).  It represents an empirical mass-shell measure,
defined transparently by its action on any smooth test function \(\psi\).
Let \(\Sigma\) be the physical hypersurface being sampled, let the same symbol
also denote the set of marker worldlines intersecting it, and let \(W_p\) be
the invariant rest-number weight of marker \(p\).  Then
\[
 \widehat f_{\Sigma}[\psi]
 :=\sum_{p\in\Sigma}W_p\,
   \psi(r_p,\pi_p,\eta_p),                              \tag{IEF45e-1a}
\]
where \(\pi_p\) denotes the chosen radial momentum coordinate:
\(\pi_p=w_p\) on an SC slice and \(\pi_p=P_p\) on an EF slice.  Equality
of the sampled distribution means equality of this measure after each marker
worldline is intersected with the same physical hypersurface and its momentum
is transformed with (IEF42)--(IEF44).  The strict certificate therefore
authenticates unique persistent identifiers (IDs) and source IDs, invariant
\(W_p\) and
\(\eta_p\), mapped event positions, and transformed momenta.  Equality of all
of those atoms is equality of (IEF45e-1a) for every test function; the code
does not pretend to reconstruct arbitrary pointwise probes from a finite
marker set.  When a crossing or trajectory table supplies both
\code{f_schwarzschild} and \code{f_eddington_finkelstein}, their scalar
values are compared too, but explicit scalar-\(f\) columns are optional.
Thus two noisy coordinate histograms are never treated as pointwise samples
of \(f\).  A separately chosen smooth reconstruction may be compared using
(IEF45d)--(IEF45e), but that is an additional diagnostic, not a replacement
for empirical-measure identity.  Stress-tensor parity is checked independently
from the deposited moments below; agreement of finitely many moments alone
would not prove equality of the distribution.

There is also a spatial Jacobian because the two distributions live on
different hypersurfaces.  The reduced canonical phase measures are
\[
 \dd\Gamma_\tau=8\pi^2e^\lambda\,
       \dd r\,\dd w\,\dd\eta,
 \qquad
 \dd\Gamma_v=8\pi^2\,
       \dd r\,\dd P\,\dd\eta .                       \tag{IEF45e-2}
\]
Hamiltonian Liouville flow maps the first measure to the second when one
follows each worldline from its \(\tau\)-slice intersection to its
\(v\)-slice intersection.  The change in the intersection radius supplies
the spatial factor that complements \(\partial P/\partial w\).  This is why
a marker's represented particle number is copied unchanged; multiplying its
weight by either momentum Jacobian would count that factor twice.

The same distinction applies to moments.  Raw coordinate components of a
tensor are not expected to have the same numbers.  In SC variables the radial
stress-energy block reconstructed from the physical moments in (EWP10) is
\[
 T^{\rm SC}_{\tau\tau}=e^{2\mu}\rho,\qquad
 T^{\rm SC}_{\tau r}=-e^{\mu+\lambda}j,\qquad
 T^{\rm SC}_{rr}=e^{2\lambda}p_{\rm rad}.             \tag{IEF45f}
\]
The sign in the middle identity agrees with
\(\lambda_\tau=-4\pi r e^{\lambda+\mu}j\).
Let \(\zeta:=A h=v_r=(ab)^{-1}\), where \(\zeta\) is the forward null
slope and the last equality follows from
(IEF38d).  Pulling the EF tensor back with
\(v_\tau=A\), \(v_r=\zeta\), and \(r_\tau=0\) gives
\begin{align}
 T^{\rm SC}_{\tau\tau}
 &=A^2T^{\rm EF}_{vv},                                      \tag{IEF45g}\\
 T^{\rm SC}_{\tau r}
 &=A\bigl(\zeta T^{\rm EF}_{vv}+T^{\rm EF}_{vr}\bigr),   \tag{IEF45h}\\
 T^{\rm SC}_{rr}
 &=\zeta^2T^{\rm EF}_{vv}+2\zeta T^{\rm EF}_{vr}
   +T^{\rm EF}_{rr}.                                      \tag{IEF45i}
\end{align}
The angular components are unchanged because the angular coordinates are
unchanged.  Equivalently, the inverse formulas follow from
\(\tau_v=F\), \(\tau_r=-F/(ab)=-h\).  These identities require all of
\(T_{vv},T_{vr},T_{rr}\), not just the two source combinations used by the
EF radial constraints.  Equations (IEF18)--(IEF20) specify how the diagnostic
deposition must construct them.
Writing the inverse formulas in terms of the polar physical moments exposes
every sign and is useful for a direct code test.  The second line below
repeats the source identity (IEF45i-2):
\begin{align}
 T^{\rm EF}_{vv}&=F^2e^{2\mu}\rho,&
 T^{\rm EF}_{vr}&=-F e^{\lambda+\mu}(\rho+j)
                 =-b(\rho+j),                        \tag{IEF45i-1}\\
 T^{\rm EF}_{rr}&=e^{2\lambda}(\rho+2j+p_{\rm rad}),&
 S&=-\frac2bT^{\rm EF}_{vr}-aT^{\rm EF}_{rr}
   =\rho-p_{\rm rad}.                                    \notag
\end{align}
For one-direction tangential pressure
\[
 p_\perp=\frac{2\pi}{r^4}\int_0^\infty\!\!\int_0^\infty
       \frac{\eta}{\epsilon}f\,
       \dd\eta\,\dd w,
 \qquad T_{\theta\theta}=r^2p_\perp,\qquad
 T_{\phi\phi}=\sin^2\theta\,T_{\theta\theta}.       \tag{IEF45i-3}
\]
Equations (IEF45e-1) and (IEF42) turn (IEF18)--(IEF23) into
(IEF45i-1)--(IEF45i-3) term by term.

The coordinate Jacobian must likewise be applied to the metric.  The pullback
of (IEF2) to \((\tau,r)\) is
\begin{align}
 (g_{\tau\tau})_{\rm pullback}&=-ab^2A^2,                 \tag{IEF45j}\\
 (g_{\tau r})_{\rm pullback}&=-ab^2A\zeta+bA,            \tag{IEF45k}\\
 (g_{rr})_{\rm pullback}&=-ab^2\zeta^2+2b\zeta.          \tag{IEF45l}
\end{align}
For an exact map these are respectively
\(-e^{2\mu},0,e^{2\lambda}\).  Thus ``the metric is the same'' means that
(IEF45j)--(IEF45l), evaluated at mapped events, agree with the SC tensor; it
does not mean that the raw arrays \((a,b)\) equal \((\lambda,\mu)\).
The areal radius, Misner--Sharp mass
\(m=r(1-g^{ab}\partial_a r\partial_b r)/2=r(1-a)/2\), and scalar contractions
of \(T_{ab}\) provide additional coordinate-independent comparisons.

\paragraph{Solving the dynamical map numerically.}
For a complete SC spacetime history, define
\(h(\tau,r)=e^{\lambda-\mu}\), impose
\(v(\tau,R_B)=\tau\) on the common outer vacuum worldtube, and solve
the advanced eikonal equation (IEF38a), equivalently
\[
 \boxed{\ v_r-hv_\tau=0,\qquad v_\tau>0\ },            \tag{IEF45m}
\]
with the plus/ingoing root.  For one ingoing null slice labelled by its
boundary value \(v_{\cal B}\), its SC time \(\tau_{v_{\cal B}}(r)\) and the
inverse Jacobian factor \(F=\partial\tau/\partial v\) may be obtained without
finite-differencing two nearly equal null labels:
\begin{align}
 \frac{\dd\tau_{v_{\cal B}}}{\dd r}
 &=-h\bigl(\tau_{v_{\cal B}}(r),r\bigr),
 &\tau_{v_{\cal B}}(R_B)&=v_{\cal B},                  \tag{IEF45n}\\
 \frac{\dd F}{\dd r}
 &=-\partial_\tau h\bigl(\tau_{v_{\cal B}}(r),r\bigr)F,
 &F(R_B)&=1.                                             \tag{IEF45o}
\end{align}
Then \(A=F^{-1}\) and \(b=F e^{\lambda+\mu}\) on that slice.  Equations
(IEF45n)--(IEF45o) are the characteristic form of (IEF38c), not an additional
assumption.

Conversely, an EF history determines the SC label from the inverse eikonal
equation
\[
 \boxed{\ \tau_r+\frac{1}{ab}\tau_v=0,\qquad \tau_v>0\ },
 \qquad \tau(v,R_B)=v .                                 \tag{IEF45p}
\]
Along a surface of constant boundary SC label \(\tau_{\cal B}\), its EF
advanced time satisfies
\[
 \frac{\dd v_{\tau_{\cal B}}}{\dd r}
 =\frac{1}{a(v_{\tau_{\cal B}}(r),r)
             b(v_{\tau_{\cal B}}(r),r)},
 \qquad v_{\tau_{\cal B}}(R_B)=\tau_{\cal B}.         \tag{IEF45q}
\]
The same-event comparator also carries the coordinate Jacobians.  Directly
advancing a Jacobian sensitivity can give a negative number when an explicit
RK stage is stiff, although the exact, future-oriented map cannot change its
sign.  The comparator therefore advances logarithmic sensitivities.  Its
forward trace starts at the sampled SC event \((\tau_*,r_*)\), not at the
outer boundary:
\begin{align}
 \frac{\dd\tau_{\rm tr}}{\dd r}&=-h(\tau_{\rm tr},r),&
 \tau_{\rm tr}(r_*)&=\tau_*,\notag\\
 \frac{\dd\ell_A}{\dd r}&=-\partial_\tau h(\tau_{\rm tr},r),&
 \ell_A(r_*)&=0,\qquad
 A(\tau_*,r_*)=e^{\ell_A(R_B)}>0 .          \tag{IEF45q-1}
\end{align}
Here the subscript ``tr'' means the value being carried along the trace.
Because \(v=\tau\) at \(R_B\), the endpoint
\(\tau_{\rm tr}(R_B)\) is the advanced-time label \(v(\tau_*,r_*)\), up to
the recorded constant time-origin offset.  The tangent equation for
\(\ell_A\) is just the derivative of the ray endpoint with respect to its
starting time.  The inverse trace starts instead at the sampled EF event
\((v_*,r_*)\):
\begin{align}
 \frac{\dd v_{\rm tr}}{\dd r}&=\zeta_{\rm inv}(v_{\rm tr},r),&
 v_{\rm tr}(r_*)&=v_*,\notag\\
 \frac{\dd\ell_F}{\dd r}&=\partial_v\zeta_{\rm inv}(v_{\rm tr},r),&
 \ell_F(r_*)&=0,\qquad
 F(v_*,r_*)=e^{\ell_F(R_B)}>0,               \tag{IEF45q-2}\\
 h&=e^{\lambda-\mu},&\zeta_{\rm inv}&:=(ab)^{-1}.      \tag{IEF45q-3}
\end{align}
At the inverse endpoint, \(v_{\rm tr}(R_B)\) gives the reconstructed SC time
\(\tau(v_*,r_*)\), again after applying the recorded time-origin offset.
Define the two terminal log Jacobians
\(q_A:=\ell_A(R_B)=\log A(\tau_*,r_*)\) and
\(q_F:=\ell_F(R_B)=\log F(v_*,r_*)\).  These are the values serialized by the
comparator.  Advancing \(\ell_A\) and \(\ell_F\) is an exact change of
variables, not an approximation or a sign repair.  A finite terminal value
always gives a positive Jacobian; a value outside floating-point range makes
the audit inconclusive and is never clipped.

The signs look opposite if one writes differential equations for the
Jacobian \emph{fields} with the outer-boundary value fixed, because the code
instead propagates an identity sensitivity from the event outward.  Along
the corresponding constant-label rays, the field equations are
\[
 \boxed{\quad
 \left.\frac{\dd\ln A}{\dd r}\right|_{v}
     =+\partial_\tau e^{\lambda-\mu},
 \qquad
 \left.\frac{\dd\ln F}{\dd r}\right|_{\tau}
     =-\partial_v\!\left(\frac1{ab}\right)
 \quad} .                                             \tag{IEF45q-3a}
\]
Thus (IEF45q-1)--(IEF45q-2) and (IEF45q-3a) describe the same derivatives
with opposite integration orientations.  Each code trace initializes its
sensitivity to one at \(r_*\), so its logarithm starts at zero.

The comparator reconstructs the \emph{same spatial metric interpolant} used
by the particle and null-map pushers.  Its comparison must therefore include
both cell-center and face data, since together they determine the metric
between stored points.  In one half-cell let \(x_0<x_1\)
be the face and center radii (or the center and face radii), set
\(\xi=(r-x_0)/(x_1-x_0)\), and attach subscripts \(0,1\) to the corresponding
endpoint data.  The production SC sampler uses
\begin{align}
 \lambda(r)&=(1-\xi)\lambda_0+\xi\lambda_1,             \tag{IEF45q-3b}\\
 \mu(r)&=(2\xi^3-3\xi^2+1)\mu_0
 +(\xi^3-2\xi^2+\xi)\Delta r\,d_0 \notag\\
 &\quad+(-2\xi^3+3\xi^2)\mu_1
 +(\xi^3-\xi^2)\Delta r\,d_1,                         \tag{IEF45q-3c}
\end{align}
where \(\Delta r=x_1-x_0\) and \(d_0,d_1\) are the saved radial derivatives
of \(\mu\).  Equation (IEF45q-3c) is an ordinary cubic joining the saved
endpoint values and slopes.  Its slope limiter is also reproduced exactly:
with secant \(\delta=(\mu_1-\mu_0)/\Delta r\), both slopes are set to zero if
\(\delta\leq0\).  Otherwise negative slopes are replaced by zero; writing
\(\alpha=d_0/\delta\) and \(\beta=d_1/\delta\), both remaining slopes are
multiplied by \(3/\sqrt{\alpha^2+\beta^2}\) when
\(\alpha^2+\beta^2>9\).  This prevents an interpolated lapse overshoot
between two valid stored points.

The code then forms
\(h=\exp[\lambda(r)-\mu(r)]\) separately on the two saved time slices and
linearly interpolates those two values in time.  The order matters:
exponentiating node values first, or linearly joining center values across a
whole cell, defines a different discrete spacetime.  On the EF side, the
production sampler linearly joins face and center values of \(a\) and \(b\)
on each half-cell; the comparator does the same before interpolating those
values in \(v\) and differentiating the resulting piecewise-linear time
history in (IEF45q-2).

The SC profile therefore saves the left/right face radii, face values of
\(\lambda,\mu,\mu_r\), and the center value of \(\mu_r\) beside every center
row.  The EF profile saves the same face locations and the left/right face
values of \(a,b\).  The authenticated \code{VENUSSCHIST1} version-2 donor
history saves the full SC face reconstruction for every negative-\(\tau\)
snapshot as well.  This last point is essential because an early same-event
ray can enter the time-reversed donor segment.  A legacy center-only profile
or version-1 donor may still be inspected diagnostically, but it is marked as
using a different center-to-center interpolant and cannot issue a production
parity certificate.  This fail-closed rule prevents a substep-independent
interpolation floor from being mistaken for a failure of RK4 convergence.

Let subscripts \(n\) and \(2n\) denote otherwise identical per-cell RK4
traces.  The history-comparison gate uses
\begin{align}
 E_{\rm fwd}
 &=\max\!\left(
      |v_{2n}-v_n|,
      |q_{A,2n}-q_{A,n}|
    \right),                                           \tag{IEF45q-4}\\
 E_{\rm inv}
 &=\max\!\left(
      |\tau_{2n}-\tau_n|,
      |q_{F,2n}-q_{F,n}|
    \right).                                           \tag{IEF45q-5}
\end{align}
Both must be at most \(2\times10^{-5}\).  No factor of \(15\) and no division
by \(|q_A|\) or \(|q_F|\) is used here.  The absolute change in \(\log A\)
directly measures the multiplicative change in \(A\); dividing it by a large
\(|\log A|\) could hide a percent-level Jacobian error.  A Richardson factor
would additionally presume a demonstrated smooth fourth-order regime, which
piecewise-linear saved histories need not have when a ray crosses a time
knot.  Increasing the substeps must make (IEF45q-4)--(IEF45q-5) converge at
the unchanged ceiling before a coordinate certificate can pass.

The implementation follows the displayed definition literally:
\codesource{Dynamic SC--EF coordinate comparator}{\path{cpu/tools/compare_wpic_dynamic_coordinates.py}}{Forward and inverse eikonal step-refinement gate}
\begin{lstlisting}[language=Python,basicstyle=\ttfamily\scriptsize]
time_change = abs(time_2n - time_n)
log_jacobian_change = abs(log_jacobian_2n - log_jacobian_n)
step_refinement_error = max(time_change, log_jacobian_change)
if step_refinement_error > eikonal_convergence_tolerance:
    status = "forward_eikonal_not_converged"  # or inverse
\end{lstlisting}
The output records the two component changes, their maximum, the substep
count, and the event with the largest error.  These values reproduce the
acceptance decision directly from the saved comparison.

For a production parity campaign, the integer \(n\) is not an implicit code
default.  The submission helper requires one explicitly chosen positive value
through \code{VENUS_EIKONAL_SUBSTEPS}.  That same value is passed to all four
resolution--realization diagnostics, written into the convergence evidence,
and required again by the terminal certificate.  Each diagnostic evaluates
both \(n\) and \(2n\), tests their raw difference in
(IEF45q-4)--(IEF45q-5), and keeps the more refined \(2n\) solution for the
same-event comparison.  A missing, zero, or inconsistent value stops the
chain; it is never replaced silently by a local default.

The calculation records the source identities for its numerical evidence.  Its Slurm
scripts and scientific Python tools are fixed under the staged
\code{VENUS_SOURCE_ROOT}.  Their SHA-256 identities are measured at
submission, checked before and after each dependent job, and carried into the
convergence record and final certificate.  Environment variables cannot
replace these production tools or Slurm scripts.  The hash checks bind each result to the comparator and convergence code
that produced it.

The two independently reconstructed maps must close:
\(\tau_{\rm EF}(v_{\rm SC}(\tau,r),r)=\tau\).  Required numerical residuals
are
\begin{align}
 R_v&=-e^{-2\mu}v_\tau^2+e^{-2\lambda}v_r^2,           \tag{IEF45r}\\
 R_\tau&=\tau_r+(ab)^{-1}\tau_v,                       \tag{IEF45s}\\
 R_{\rm inv}&=\tau_{\rm EF}(v_{\rm SC}(\tau,r),r)-\tau.\tag{IEF45t}
\end{align}
They must be reported, rather than assuming the static tortoise formula in a
time-dependent region.
The Jacobian part of the closure test is evaluated without subtracting two
large products,
\[
R_{AF}=\left|\operatorname{expm1}(q_A+q_F)\right|
       =|AF-1|,                                        \tag{IEF45t-1}
\]
Here \(\operatorname{expm1}(x)\) is the floating-point operation that
evaluates \(e^x-1\) accurately when \(x\) is small.
Any exponential overflow fails closed.  The comparator's dimensionless
time-closure gate is exactly
\[
 E_{\rm comp}^{\tau}
 =\frac{|\tau_{\rm EF}(v_{\rm SC}(\tau,r),r)-\tau|}
        {\max(1,|\tau|,|\tau_{\rm EF}|)},
 \qquad
 \max(E_{\rm comp}^{\tau},R_{AF})\le5\times10^{-4}    \tag{IEF45t-2}
\]
by default.  Thus the raw eikonal step-refinement test
(IEF45q-4)--(IEF45q-5) and the normalized map-composition test are distinct
quantities and are reported separately.

\paragraph{Dynamical initial data on an ingoing null surface.}
The coordinate-neutral marker file currently used by the driver fixes the
same random realization, weights, \(r,w,\eta\), and IDs.  That is valuable
for controlling particle-in-cell (PIC) noise, but by itself it is
\emph{not} a transformed
dynamical initial data set: inserting it on both \(\tau=0\) and \(v=0\)
places the markers at different spacetime events.  A production comparison of the same physical evolution requires
characteristic extraction:
\begin{enumerate}
\item Evolve and densely checkpoint the SC metric and every marker worldline
      over the time interval crossed by a chosen ingoing null surface.
\item Solve (IEF45n)--(IEF45o) from the outer vacuum worldtube.  For every
      marker ID, find the unique intersection of its timelike worldline with
      \(\tau=\tau_{v_{\cal B}}(r)\).  A bracketed scalar root may be used:
      \(G_p(\tau)=v(\tau,r_p(\tau))-v_{\cal B}\).  Since the SC
      characteristic has
      \(\dd r_p/\dd\tau=e^{\mu-\lambda}w_p/\epsilon_p\),
      \[
       \frac{\dd G_p}{\dd\tau}
       =A\left(1+\frac{w_p}{\epsilon_p}\right)>0,     \tag{IEF45u}
      \]
      where \(\epsilon_p>|w_p|\) for a massive marker.  Hence a crossing,
      once bracketed, is unique and future oriented.
\item At each intersection retain its invariant rest weight, \(\eta\), and
      IDs, and transform its covector using (IEF42).  Interpolate the SC metric
      to the same null surface and construct \(a=e^{-2\lambda}\) and
      \(b=F e^{\lambda+\mu}\).
\item Deposit the extracted marker ensemble and require the EF constraint
      solution to converge to those transformed metric and stress-tensor
      data.  Only this extracted artifact is a valid EF initial state for the
      dynamical parity run.
\item In the common future, solve (IEF45m) and (IEF45p), intersect paired
      marker worldlines with common physical slices, and compare marker
      \((r,w,\eta)\), the scalar distribution via (IEF45d), all moments via
      (IEF45f)--(IEF45i), and the metric via (IEF45j)--(IEF45l).
\end{enumerate}
The SC history must extend far enough that the entire null slice, including
its innermost point, lies in the stored domain.  A surface truncated before
all marker intersections is not valid characteristic data.

\paragraph{Well-posedness and discretization domain.}
On a smooth polar exterior, \(c=e^{\mu-\lambda}\) is positive and locally
Lipschitz.  The ray equation in (IEF38c-1) then has unique solutions, so two
radial ingoing characteristics cannot cross at \(r>0\).  The spherical
congruence does focus at \(r=0\), which is its ordinary central caustic and
the endpoint of this one-sided areal-radius map.  The other endpoint is a
polar-coordinate horizon: the construction requires
\[
 a=e^{-2\lambda}>0,\qquad A=v_\tau>0,\qquad
 \det\frac{\partial(v,r)}{\partial(\tau,r)}=A,        \tag{IEF45v}
\]
and becomes singular in the polar variables as \(a\to0^+\).  These are
limitations of the overlap chart, not of the EF evolution.  Eventwise parity
is required everywhere that (IEF45v) holds; after the EF run crosses
\(a=0\), there is no SC chart against which to compare it.

For a stored SC history, VENUS may integrate (IEF45n) with an adaptive
ordinary differential equation (ODE)
scheme and spacetime interpolation of \(h\).  Integrating
\[
 \frac{\dd\log F}{\dd r}=-\partial_\tau h,\qquad
 \log F(R_B)=0                                           \tag{IEF45w}
\]
instead of \(F\) preserves \(F>0\) exactly.  An equivalent full-grid method
evolves \(s=v_r>0\) with the conservative equation
\[
 s_\tau-\partial_r(cs)=0,\qquad A=cs,               \tag{IEF45x}
\]
using a positivity-preserving upwind flux.  In either formulation, the
spacetime interpolation, marker dense output, and scalar root solve must be
at least as accurate as the SC and EF discretizations being compared.
Otherwise the coordinate mapper, rather than either Einstein--Vlasov solver,
sets the observed parity error.

No coverage percentage by itself establishes coordinate equivalence.  The
implemented minimum certificate is a \(2\times2\) matrix: two jointly refined
radial/time resolutions and two independent deterministic quiet-PIC
realizations at each resolution.  The resolution pair estimates truncation
change, while the realization pair exposes the sampled-measure floor; neither
pair alone substitutes for the other.  A third resolution or additional
realizations are strongly recommended near marginal tolerances, but are not
silently claimed by the current automated gate.  The raw eikonal
step-refinement changes and the inverse-closure errors must be smaller than
the field interpolation error;
the SC--EF differences in pulled-back metric, stress tensor,
Misner--Sharp mass, and every available marker audit must decrease under
refinement and become no larger
than the authenticated truncation envelope plus the measured PIC floor.
Cells with negligible matter are excluded from relative moment norms but
retained in absolute and metric norms.  Passing only fingerprints, ADM mass,
compactness milestones, one resolution, or one realization is a smoke test,
not proof of dynamical coordinate parity.

Here a PIC \emph{realization} means one deterministic choice of the
Cranley--Patterson shifts in (IEF-S23).  It samples the same continuum family
with different marker locations and weights; it is not a different physical
model.  A \emph{resolution level} means one joint choice of radial spacing,
maximum timestep, and marker count.  Coordinate convergence means that the
same-event map, metric, stress/moment, and available empirical-measure errors
decrease from the coarse level to the fine level until the change is no
larger than the difference between the two realizations.  This explains why
neither two grids with one realization nor two realizations on one grid are
sufficient.

\paragraph{Exact common-event domain.}
The coordinate comparison covers the connected exterior in which both
histories describe the same physical events.  Physical chart and data
boundaries determine this domain.  A numerical map failure within it remains
a failed comparison.  Let
\begin{align}
 r_-&=\max(r_{\min}^{\rm SC},r_{\min}^{\rm EF},r_{\min}^{\rm request}),&
 r_+&=\min(r_{\max}^{\rm SC},r_{\max}^{\rm EF},R_B,
           r_{\max}^{\rm request}).                    \tag{IEF45y}
\end{align}
The code samples ordered SC times \(\{\tau_i\}\) and ordered radii
\(\{r_j\}\subset[r_-,r_+]\).  At a fixed \(\tau_i\), it scans inward from the
largest sampled radius.  Only three physical or data boundaries may end the
candidate exterior: the complete forward ray has left the stored SC history,
the mapped event has left the stored EF history, or the overlap chart has
reached
\[
 a_{\rm SC}(\tau,r)>a_{\min},\qquad
 a_{\rm EF}(v(\tau,r),r)>a_{\min},\qquad
 a_{\min}=10^{-3}\ \hbox{by default}.                  \tag{IEF45z}
\]
These boundaries have the machine-readable statuses
\path{outside_sc_eikonal_domain}, \path{outside_ef_domain}, and
\path{at_or_beyond_horizon}.  The radial candidate set \(D_i\) extends inward from the common outer
worldtube to the first such boundary.  It therefore forms one unbroken
exterior interval, fixed before the numerical errors are assessed.

A failed integrator, failed inverse composition, or failed metric identity is
\emph{not} an allowed boundary.  Such a sample remains inside \(D_i\) as an
unexplained hole and fails the certificate.  In particular, an event whose
raw forward step-doubling change (IEF45q-4) is too large keeps the status
\path{forward_eikonal_not_converged}; it is never re-labelled as outside the
EF history by using the time returned by that nonconverged trace.  Likewise,
valid-looking samples that reappear inward of a non-horizon boundary are
reported as disconnected support rather than being silently admitted.

Let \(N_r\) be the number of sampled radii, and let \({\cal T}\) be the
longest contiguous block of sampled times for which each exterior suffix
covers at least 95 percent of those radii.  A tie is resolved in favour of
the earlier block.  With \(N_i=|D_i|\), \(N_i^{\rm map}\) the number of
successfully mapped events in that suffix, and \(N_{\rm full}\) the number
of time slices with every event in their suffix mapped, the reported
coverages are exactly
\begin{align}
 C_r&=\min_{i\in{\cal T}}\frac{N_i}{N_r},&
 C_{\rm event}&=
  \frac{\sum_{i\in{\cal T}}N_i^{\rm map}}
       {\sum_{i\in{\cal T}}N_i},&
 C_{\rm time}&=\frac{N_{\rm full}}{|{\cal T}|}.       \tag{IEF45y-1}
\end{align}
The strict defaults require
\[
 C_r\ge0.95,\qquad C_{\rm event}\ge0.95,\qquad
 C_{\rm time}\ge0.95,\qquad
 \sum_{i\in{\cal T}}N_i^{\rm map}\ge64,\qquad
 N_{\rm full}\ge3,                                   \tag{IEF45y-2}
\]
\emph{and} zero unexplained holes.  The apparently redundant coverage
fractions are retained in the certificate so that the sampled scope remains
explicit if the sampling policy is extended later.  A gap followed by a
later qualifying time slice, re-entry after a non-horizon radial boundary, or
loss of common support before a chart/horizon endpoint all fail closed.
Initial slices outside the overlap of the two outer clocks and terminal slices
after a genuine chart cutoff are counted and reported, but are not presented
as compared events.

Only after \(D_i\) and \({\cal T}\) have been fixed does the comparator rebuild
the metric, stress, moment, and marker accumulators.  Consequently every norm
contains exactly the selected common-domain events.  The maximum forward and
inverse eikonal changes and the inverse-composition errors used as gates are
also reduced only on this declared domain.  Separate
\path{all_attempted_maximum_*} fields and the \path{excluded} counters retain
failures outside it, including outer-clock mismatch, eikonal-domain exits,
forward or inverse nonconvergence, inverse-composition failure, the horizon
cutoff, an SC metric-identity failure, and successfully mapped events outside
the chosen maximal block.  This separation keeps errors within the declared comparison domain in its
norms and records all excluded attempts separately.

The final certified common time is therefore the end of \({\cal T}\), subject
to all of (IEF45y-2) and the no-hole/topology tests.  It is never extended by
copying the EF clock after the polar chart reaches its horizon barrier.
Events with \(a\le a_{\min}\), and all post-horizon EF events, are explicitly
outside the coordinate certificate.

\paragraph{Saved-output cadence.}
The eikonal map interpolates between saved times, so output cadence is part of
the numerical evidence.  The strict default requirement is
\[
 \frac{\max\Delta\tau_{\rm save}}{\min\Delta r_{\rm SC}}\le0.75,
 \qquad
 \frac{\max\Delta v_{\rm save}}{\min\Delta r_{\rm EF}}\le0.75. \tag{IEF45aa}
\]
If every \(k\)-th accepted step is saved, the conservative setup check
\(k\,\Delta t_{\max}/\min\Delta r\le0.75\) is sufficient.  A sparse
science-output cadence can resolve horizon/halo trends while still failing
(IEF45aa); such output must never be reused as a strict coordinate
certificate.

\paragraph{Paired-coordinate visualization and its provenance.}
The tool
\path{cpu/tools/render_weighted_einstein_vlasov_coordinate_pair_gif.py}
builds the standard paired Graphics Interchange Format (GIF) animation.  It
does not align equal printed clock values.  For every displayed overlap frame
it authenticates the solved-null manifest and donor-v2 history, solves and
checks both the forward and inverse event maps, and samples SC and EF curves
at the same spacetime events.  Its three panels show
\(\log_{10}\rho\) against \(\log_{10}r\), the logarithmic density slope
against \(\log_{10}r\), and compactness \(2m/r\) against linear \(r\), where
\(\log_{10}\) is the base-10 logarithm.

After the SC chart ends, the renderer may continue with saved EF slices, but
those frames are clearly labelled EF-only and are not presented as coordinate
comparisons.  It writes Portable Network Graphics (PNG) preview frames, an
exact frame-clock sidecar, and SHA-256 provenance sidecars for the inputs and
animation.  Therefore the GIF is a reproducible visualization of the
same-event audit and post-horizon continuation; it is not itself a substitute
for the strict coordinate certificate.

\paragraph{Output cadence for coordinate comparisons.}
The coordinate comparison requires output dense enough to resolve the null
map. For a linear field grid with \(R_E=24\), \(N_r=2048\), and
\(\Delta t_{\max}=0.004\), the cell width is
\(\Delta r=24/2048=0.01171875\). Equation (IEF45aa) then permits
at most \(\code{high-output-steps}=2\). Saving every step satisfies
this condition. A coarser pair with \(N_r=1024\) and
\(\Delta t_{\max}=0.008\) has the same ratio. An outcome history may
use a sparser cadence, while the coordinate comparison must satisfy its
own sampling criterion.

\paragraph{What is and is not proved about the evolving distribution.}
The required, hashed marker-crossing table proves empirical-measure equality
on the imported initial null surface.  Later same-event metric and all
stress/moment components are tested throughout the common domain.  The
certifier can additionally audit an authenticated
\path{sampled_marker_trajectories} artifact on later mapped null surfaces, but
the present production writer does not emit that optional artifact.  It
therefore records
\path{later_trajectory_evidence_present=false}; a passing current certificate
must not be described as marker-by-marker or pointwise-\(f\) equality at later
times.  Obtaining that stronger numerical statement requires implementing and
hashing the later trajectory table.  This limitation does not weaken the
later metric and moment comparisons that are actually performed.

\subsubsection{Finite-volume WPIC deposition}

The module stores a distinct marker record.  For example, an active marker is
constructed through the public type and exact field names as follows:
\codesource{EF marker data object}{\path{cpu/include/ev/weighted_einstein_vlasov_pic_eddington_finkelstein.hpp}}{Public marker fields, state, and validity convention}
\begin{lstlisting}[language=C++,basicstyle=\ttfamily\scriptsize]
namespace ef = ev::eddington_finkelstein;

ef::Marker marker;
marker.rest_weight = represented_rest_weight;
marker.r = initial_areal_radius;
marker.p_r = initial_canonical_p_r; // strictly positive
marker.eta = initial_L_squared / 2.0;
marker.id = persistent_id;
marker.source_id = source_cell_id;
marker.state = ef::MarkerState::active;
\end{lstlisting}
The conservative radial mesh uses the volume coordinate
\[
 \xi=\frac{r^3}{3},\qquad
 \Delta V_i=4\pi\Delta \xi_i
 =\frac{4\pi}{3}(r_{i+1}^3-r_i^3).                    \tag{IEF46}
\]
The corresponding public constructor is
\codesource{EF finite-volume field grid}{\path{cpu/include/ev/weighted_einstein_vlasov_pic_eddington_finkelstein.hpp}}{Public grid constructor and origin-parity mesh convention}
\begin{lstlisting}[language=C++,basicstyle=\ttfamily\scriptsize]
namespace ef = ev::eddington_finkelstein;
ev::WeightedPICFieldGrid grid =
    ef::make_field_grid(params, radial_cells);
\end{lstlisting}
Let \(\omega_{pi}\) be the CIC or TSC cloud fraction deposited by active
marker \(p\) into cell \(i\).  CIC (cloud in cell) is the two-cell linear
shape, and TSC (triangular shaped cloud) is the smoother three-cell quadratic
shape.  The fractions are nonnegative and sum to one over the receiving
cells for a marker whose stencil is wholly on the mesh:
\(\sum_i\omega_{pi}=1\).  Then the implemented shell-extensive sources
are
\begin{align}
 {\cal W}_i&=\sum_p\omega_{pi}W_p,                     \tag{IEF47}\\
 {\cal R}_i&=\sum_p\omega_{pi}W_pP_p,                  \tag{IEF48}\\
 {\cal S}_i&=\sum_p\omega_{pi}W_p
       \frac{1+2\eta_p/r_p^2}{P_p}.                    \tag{IEF49}
\end{align}
The cell sources are
\[
 (T_{rr})_i=\frac{{\cal R}_i}{\Delta V_i},\qquad
 S_i=\frac{{\cal S}_i}{\Delta V_i}.                   \tag{IEF50}
\]
Unequal-weight sampling quality is reported as
\[
 N_{{\rm eff},i}
 =\frac{(\sum_p\omega_{pi}W_p)^2}
        {\sum_p(\omega_{pi}W_p)^2}.                    \tag{IEF51}
\]
Excised and escaped records remain in rest-weight diagnostics but do not
deposit into \(T_{rr}\) or \(S\).  Parallel workers construct marker stencils,
then the code accumulates them in deterministic marker-array order using
long-double partial sums.  Consequently source arrays are independent of the
requested thread count for a fixed marker ordering.

The source application programming interface (API) is
\codesource{EF WPIC source deposition}{\path{cpu/include/ev/weighted_einstein_vlasov_pic_eddington_finkelstein.hpp}}{Deterministic CIC/TSC deposition into \(T_{rr}\) and \(S\)}
\begin{lstlisting}[language=C++,basicstyle=\ttfamily\scriptsize]
namespace ef = ev::eddington_finkelstein;

ef::DepositedSources sources = ef::deposit_sources_parallel(
    markers,
    grid.mesh,
    ev::WeightedPICParticleShape::tsc,
    threads);

// sources.t_rr[i] and sources.s[i] are the cell-volume densities.
// sources.shell_t_rr_numerator and shell_s_numerator are extensive.
\end{lstlisting}

\subsubsection{Finite-volume radial constraint integration}

Within cell \(i\), VENUS holds \(T_{rr,i}\) and \(S_i\) constant.  Equation
(IEF28) is then analytic:
\[
 b(r)=b(r_i)\exp\left[
 2\pi T_{rr,i}(r^2-r_i^2)\right],
 \qquad r_i\le r\le r_{i+1}.                           \tag{IEF52}
\]
Starting from \(b(r_{\max})=1\), face values are integrated inward as
\[
 b_i=b_{i+1}\exp\left[-2\pi T_{rr,i}
             (r_{i+1}^2-r_i^2)\right].                 \tag{IEF53}
\]
This exponential representation makes \(b>0\) manifest, including inside
\(a=0\).

The \(Y=mb\) source contains the exponential \(b(r)\).  VENUS integrates each
left and right center--face segment with deterministic three-point
Gauss--Legendre quadrature.  It first computes \(b\) from the outer boundary,
then either integrates \(Y\) outward from zero on the regular first slice or
inward from fixed \(M_{\rm ADM}\) thereafter.  Center and face values of
\(m=Y/b\), \(a=1-2m/r\), \(a_r\), and \(b_r\) are then formed from
(IEF5) and (IEF30).

For one segment \([r_L,r_R]\) inside cell \(i\), let
\(r_m=(r_L+r_R)/2\) and \(d=(r_R-r_L)/2\).  With nodes
\(z_q=\{-\sqrt{3/5},0,\sqrt{3/5}\}\) and weights
\(w_q=\{5/9,8/9,5/9\}\), the increment is
\[
 \Delta Y\simeq d\sum_{q=1}^3w_q\,
 2\pi(r_m+dz_q)^2 b(r_m+dz_q)(T_{rr,i}+S_i).
                                                        \tag{IEF53a}
\]
This quadrature integrates the exponential \(b\) from (IEF52), so it
uses the same within-cell source as the analytic \(b\) solve.
\codesource{EF radial mass-source quadrature}
{\path{cpu/include/ev/weighted_einstein_vlasov_pic_eddington_finkelstein.hpp}}
{Evaluates the three-point Gauss--Legendre contribution to \(Y=mb\) on one center--face segment}
\begin{lstlisting}[language=C++,basicstyle=\ttfamily\scriptsize]
const double r = midpoint + half_width * node[q];
const double b = b_inside_cell(
    r, cell_left, cell_left_b, t_rr);
integral += static_cast<long double>(weight[q])
    * (2.0L * static_cast<long double>(PI))
    * static_cast<long double>(t_rr + s)
    * static_cast<long double>(r) * r
    * static_cast<long double>(b);
\end{lstlisting}
The function returns \code{half_width * integral}; its variables
\code{t_rr} and \code{s} are the constant cell sources.

The public calls make the two boundary modes explicit:
\codesource{EF radial constraint solver}{\path{cpu/include/ev/weighted_einstein_vlasov_pic_eddington_finkelstein.hpp}}{Reusable regular-center and fixed-ADM-mass library modes}
\begin{lstlisting}[language=C++,basicstyle=\ttfamily\scriptsize]
namespace ef = ev::eddington_finkelstein;

// Regular first slice: infer ADM mass with the outward mb solve.
ef::ConstraintOptions first_constraints;
ef::Field first = ef::deposit_and_solve_field(
    markers, grid, ev::WeightedPICParticleShape::tsc,
    threads, first_constraints);
const double conserved_M = first.metric.adm_mass;

// Every continuation slice: fixed-M inward solve.
ef::ConstraintOptions continuation;
continuation.fixed_adm_mass = conserved_M;
continuation.excision_a_threshold = -0.01;
continuation.excision_buffer_cells = 4;
continuation.excision_horizon_fraction = 0.1;
ef::Field field = ef::deposit_and_solve_field(
    markers, grid, ev::WeightedPICParticleShape::tsc,
    threads, continuation);
\end{lstlisting}

For \(r\ge r_{\max}\), \code{sample_metric} returns the exact neutral vacuum
continuation
\[
 a=1-\frac{2M_{\rm ADM}}r,\qquad b=1,\qquad
 a_r=\frac{2M_{\rm ADM}}{r^2},\qquad b_r=0.             \tag{IEF54}
\]

\subsubsection{Interpolation, pusher, and self-consistent time coupling}

The metric object stores center and face arrays for \(a,b,a_r,b_r,m\).  Inside
a cell, \code{sample_metric} uses a piecewise-linear face-to-center or
center-to-face interpolation, never an interpolation through a factor
\(1/a\).  Outside the outer face it uses (IEF54).  The characteristic API is
\codesource{EF metric sampler and characteristic RHS}{\path{cpu/include/ev/weighted_einstein_vlasov_pic_eddington_finkelstein.hpp}}{Face-aware metric interpolation and equations (IEF15)--(IEF16)}
\begin{lstlisting}[language=C++,basicstyle=\ttfamily\scriptsize]
namespace ef = ev::eddington_finkelstein;

ef::CharacteristicState state{marker.r, marker.p_r};
ef::MetricSample local = ef::sample_metric(state.r, grid, field.metric);
ef::CharacteristicDerivative rhs = ef::characteristic_rhs(
    state, marker.eta, local);
\end{lstlisting}

The public \code{CharacteristicState} continues to store \(P=p_r\), and
\code{characteristic_rhs} returns \((\dot r,\dot P)\) from
(IEF15)--(IEF16).  Internally, however, \code{rk4_advance} pushes
\[
 u:=\log P,\qquad P=e^u,\qquad
 \dot r=\frac b2\left(a-\chi e^{-2u}\right),          \tag{IEF54a}
\]
and the exact quotient \(\dot u=\dot P/P\):
\[
 \boxed{
 \dot u=-\frac12(ab_r+ba_r)
        -\frac12 b_r\chi e^{-2u}
        +\frac{2b\eta}{r^3}e^{-2u}.
 }                                                     \tag{IEF54b}
\]
Equations (IEF54a)--(IEF54b) are algebraically identical to the continuum
characteristics wherever \(P>0\); the logarithm is an implementation
variable, not a change of physics or of the public marker format.

Direct evaluation of the logarithmic right-hand side keeps the momentum
factors in the algebraically reduced expression:
\codesource{EF logarithmic-momentum characteristic}
{\path{cpu/include/ev/weighted_einstein_vlasov_pic_eddington_finkelstein.hpp}}
{Computes the exact quotient \(\dot P/P\) in the internal RK4 variable}
\begin{lstlisting}[language=C++,basicstyle=\ttfamily\scriptsize]
result.dlog_p_r =
    -0.5 * (sampled.a * sampled.b_r
        + sampled.b * sampled.a_r)
    -0.5 * sampled.b_r * chi * inverse_p_squared
    + 2.0 * sampled.b * eta * inverse_p_squared
        / (radius * radius * radius);
\end{lstlisting}
Here \code{inverse_p_squared} is \(P^{-2}=e^{-2u}\) and
\code{sampled} is the local metric from the current frozen field.

Classical fourth-order Runge--Kutta is applied to \((r,u)\), optionally with
multiple equal characteristic substeps, and \(P\) is reconstructed with
\(\exp u\).  Every stage therefore remains on the future-directed \(P>0\)
branch by construction, so an explicit stage cannot overshoot through the
singular boundary \(P=0\).  The code rejects a stage if \(r<0\), if \(u\) is
nonfinite or outside the logarithmic interval bounded by the smallest
positive normal and largest finite \code{double}, or if either transformed
right-hand side is nonfinite.
This is a hard rejection; \code{rk4_advance} does not adaptively retry with a
smaller substep.  The production driver catches a retryable characteristic,
RK4, or logarithmic-state rejection at the self-consistent-step boundary,
halves the whole \(\Delta v\), and retries at most 24 times. An explicit
inner-field-boundary violation with excision enabled terminates the evolution.
In the interior-evolution
mode, inner-boundary and recognized constraint failures enter the same
bounded retry procedure described in Section~\ref{sec:ef-interior-evolution}.

The logarithmic variable preserves the momentum sign.  Accuracy still
requires resolving the factors \(e^{-2u}=P^{-2}\), which become stiff as
\(P\) becomes small.  Timestep, substep, and radial-resolution comparisons
must therefore resolve this coordinate singularity.  Likewise, RK4 in \(u\) is not a
symplectic integrator and is numerically different from applying RK4
directly to \(P\), though both are fourth order for a smooth resolved
characteristic.  The momentum timescale in (IEF55),
\(P/|\dot P|=1/|\dot u|\), remains the relevant timestep control.

\code{advance_frozen_metric_parallel} advances active markers in parallel
while checking that rest weight, \(\eta\), ID, and source ID are unchanged.
Each batch validates the radial mesh and its agreement with the metric's
cell centers and faces before advancing any marker. A \code{MetricSampler}
then reuses this fixed mesh and metric for all particles and Runge--Kutta
stages. The standalone sampling and characteristic routines also validate
their inputs at entry. Reusing the validated sampler keeps the mesh checks
outside the particle loop:
\codesource{Sampling a fixed EF metric within a particle batch}
{\path{cpu/include/ev/weighted_einstein_vlasov_pic_eddington_finkelstein.hpp}}
{Validates the mesh once and retains the same center/face interpolation}
\begin{lstlisting}[language=C++,basicstyle=\ttfamily\scriptsize]
const ef::MetricSampler sample(grid, metric);
const auto local = sample(marker.r);
\end{lstlisting}
The referenced grid and metric remain unchanged throughout that batch.

One self-consistent interval is a midpoint field coupling:
\begin{enumerate}
\item deposit the old markers and solve the old field;
\item push a copy by \(\Delta v/2\) in the old field;
\item deposit that copy and solve the midpoint field with fixed ADM mass;
\item push the original markers by \(\Delta v\) in the midpoint field;
\item deposit the endpoint and solve the endpoint field with fixed ADM mass;
\item apply any newly available causal excision and recompute the endpoint
      field without the excised marker sources.
\end{enumerate}
The frozen characteristic is fourth order, but the coupled particle--field
update is second order in time.

The reusable library can infer the ADM mass on a regular first slice; that
result must seed every later library call.  This example exercises that
library mode.  A fresh production EF run instead sets
\code{options.fixed_adm_mass=M_null} before its first call.
\codesource{EF self-consistent library step}{\path{cpu/include/ev/weighted_einstein_vlasov_pic_eddington_finkelstein.hpp}}{Midpoint particle--field coupling; production supplies \(M_{\rm null}\) immediately}
\begin{lstlisting}[language=C++,basicstyle=\ttfamily\scriptsize]
namespace ef = ev::eddington_finkelstein;

ef::EvolutionOptions options;
options.shape = ev::WeightedPICParticleShape::tsc;
options.threads = threads;
options.characteristic_substeps = 2; // production value
options.enable_excision = true;
options.excision_a_threshold = -0.01;
options.excision_buffer_cells = 4;
options.excision_horizon_fraction = 0.1;

// options.fixed_adm_mass is empty only for this regular first call.
ef::StepResult step = ef::advance_self_consistent(
    markers, dt, grid, options);
options.fixed_adm_mass = step.conserved_adm_mass;

// Reuse options, now carrying fixed M, on every later call.
step = ef::advance_self_consistent(markers, next_dt, grid, options);
\end{lstlisting}

\subsubsection{Timestep selection}

The public estimator \code{choose_timestep} examines every active marker.
For the containing cell width \(\Delta r_i\), it limits the requested
\(\Delta v_{\max}\) by
\[
 \Delta v\le C_{\rm CFL}\frac{\Delta r_i}{|\dot r|},
 \qquad
 \Delta v\le C_{\rm CFL}\frac{P}{|\dot P|},            \tag{IEF55}
\]
and by the additional conservative acceleration heuristic
\[
 \Delta v\le
 \sqrt{\frac{2C_{\rm CFL}\Delta r_i}{|\dot P|}}.       \tag{IEF56}
\]
Zero denominators simply omit the corresponding bound.  The implementation
requires \(0<C_{\rm CFL}\le1\) and a positive finite maximum step.
\(C_{\rm CFL}\) is the Courant--Friedrichs--Lewy safety factor: it limits how
large a fraction of a cell or local momentum-change time one step may cover.
Before an excision surface is applied, an inward-moving marker also obeys
\[
 \Delta v\le C_{\rm CFL}\frac{r-r_{\rm in}}{|\dot r|},
 \qquad \dot r<0,                                      \tag{IEF56a}
\]
where \(r_{\rm in}\) is the inner field face and dots denote derivatives
with respect to advanced time \(v\). This bound resolves the distance to
the inner boundary even within the first cell. Equations (IEF55)--(IEF56a)
set the numerical timestep.

\codesource{EF timestep estimator}{\path{cpu/include/ev/weighted_einstein_vlasov_pic_eddington_finkelstein.hpp}}{Public CFL, momentum-timescale, and acceleration bounds}
\begin{lstlisting}[language=C++,basicstyle=\ttfamily\scriptsize]
const double dt = ef::choose_timestep(
    markers, grid, field.metric, 0.20, dt_max);
\end{lstlisting}
The production driver uses \(C_{\rm CFL}=0.20\) and two characteristic
substeps.  These are production controls, not universal library constants.
The resolution study must vary the CFL factor, maximum step, and RK4
substeps independently of radial and marker resolution.

\subsubsection{Apparent horizon, excision, and origin handling}

The metric solver searches all face intervals for a negative-to-positive
zero of \(a\), retaining the outermost one.  It linearly interpolates
\(r^2a\) between the two faces and records
\code{apparent_horizon_radius} and
\[
 m_H=\frac{r_H}{2}.                                     \tag{IEF57}
\]
This is an apparent-horizon diagnostic based on equation (IEF8), not an event
horizon finder.

Once an outer cell center satisfies
\[
 a<a_{\rm ex}<0,                                        \tag{IEF58}
\]
the implementation searches inward for a cell center with a fully trapped
outward buffer and radius
\[
 r_{\rm ex}\le F r_H,\qquad 0<F\le1.                    \tag{IEF-I1}
\]
Here \(F\) is the maximum excision-radius fraction selected by
\code{--ef-excision-horizon-fraction}; its default is \(0.1\).
The other defaults are \(a_{\rm ex}=-0.01\) and four outward buffer cells.
Every center and intervening face from the candidate boundary through the
buffer must satisfy (IEF58), with finite positive \(b\) at each knot.
The metric sampler interpolates linearly between these knots, so the same
conditions hold throughout the buffer.  The near-horizon choice is obtained with \(F=1\), while
the trapped-buffer condition applies to every value.  If the cap lies in an
untrapped inner cavity, excision waits for a qualifying trapped center.
The mesh must resolve the radius cap and its buffer.

Removing markers changes the stress deposited in the endpoint constraint
solve.  The retained boundary is therefore checked again in that recomputed
metric.  This includes a boundary used during the midpoint push and one
introduced by endpoint excision.  The larger retained radius must obey the
trapped-buffer condition and the cap \(F r_H\).  An unsuccessful check
rejects the step before committing the marker state.
\codesource{Recheck the retained excision surface}
{\path{cpu/include/ev/weighted_einstein_vlasov_pic_eddington_finkelstein.hpp}}
{Validate the actual endpoint metric after removing excised sources}
\begin{lstlisting}[language=C++]
retain_causal_excision_boundary(
    result.endpoint_field.metric, applied_excision_radius,
    fixed_constraints);
\end{lstlisting}

Only active markers at or inside the selected surface change to
\code{MarkerState::excised}. Their records and invariant rest weights remain
in the accounting. The active exterior sources and trajectories continue to
evolve, and the fixed outer mass in (IEF36) retains the gravitational mass
inside the removed region. The stored metric extension inside that boundary
serves the numerical scheme; physical interior plots use the retained cells.

Because \(a<0\) makes both future radial null rays move toward decreasing
areal radius, an excision boundary placed strictly inside the trapped region
is causal: no physical inner boundary data should be supplied.  Excision must
not be triggered from a merely small positive \(a\), and the result must be
checked under changes of \(a_{\rm ex}\), buffer width, radius fraction
\(F\), and radial resolution. Sensitivity comparisons hold the other
controls and physical initial data fixed.

The EF outer boundary uses \code{outflow_vacuum}. The production driver
sets \path{require_vacuum_outer_boundary=true}: a trial in which a
positive-weight marker escapes rejects the entire step and saves the last
accepted state. The summary records
\code{ef_outer_boundary_evolution_limit}, together with the final time
and \code{vacuum_outer_boundary_limit} as its cause. This applies with
excision enabled or disabled. The requested reflecting EF wall is rejected
when the field grid is constructed.

This domain condition follows from the fixed outer mass and normalization.
Once matter crosses the finite outer face, its enclosed mass and the exterior
stress must be evolved consistently. Increasing the domain keeps the isolated
vacuum boundary applicable. The reusable frozen-field particle routines can
record escaped markers and retain their weights for bookkeeping. A library
caller can select that behavior with
\path{require_vacuum_outer_boundary=false}; a physical self-consistent
outflow calculation also requires a model for the exterior field.

The ingoing-null coordinate chart is delicate at the true spherical origin.
Andr\'easson and Rein note that its naive Cartesian form is not regular there.
The present implementation starts with shell support away from \(r=0\).
Excision, when enabled, removes interior markers only after the local trapped
region satisfies the conditions above.  When \code{using_origin_parity_inner(params)} is true,
\code{make_field_grid} deliberately constructs a linear mesh with a true
\(r=0\) face even if the configured phase space \code{r_min} is positive;
otherwise it retains the configured finite-radius inner face.  This is a mesh
convention. EF characteristics follow their ingoing-null equations through
positive radii. An attempted inner crossing rejects the trial state, and
\(r=0\) is singular for \(\eta>0\). The Schwarzschild module uses a
separate Cartesian continuation at the origin.

\paragraph{Retaining matter inside the horizon.}
\label{sec:ef-interior-evolution}
The characteristic equations remain finite at \(a=0\) for finite positive
\(P\) and \(b\). A horizon therefore permits continued evolution of the
interior matter. With \code{--ef-evolve-interior}, every old, midpoint, and
endpoint field solve retains the active markers and sets its applied
excision radius to zero. The library controls are
\codesource{Retain the interior markers}
{\path{cpu/include/ev/weighted_einstein_vlasov_pic_eddington_finkelstein.hpp}}
{Select the same mode in standalone constraints and self-consistent evolution}
\begin{lstlisting}[language=C++]
ef::ConstraintOptions constraints;
constraints.enable_excision = false;
ef::EvolutionOptions evolution;
evolution.enable_excision = false;
\end{lstlisting}
The evolution routine propagates its setting to each constraint solve.
An active marker may move below the first cell center, subject to (IEF56a).
Resolving its trajectory still requires resolving the stress and metric that
determine that trajectory. In particular, the first-cell metric interpolates
the regular-origin face value \(a(0)=1\). Matter concentrated within that
cell requires radial refinement before its density profile can be interpreted.

The driver retries a failed characteristic, inner crossing, or recognized
nonfinite constraint solve with half the timestep, for at most 24 halvings.
If the trial still fails, or advanced time cannot increase in floating-point
arithmetic, it saves the complete last accepted state. The summary records
\path{interior_evolution_limit}, its cause, the final accepted time,
and the requested end time. Use \code{--expected-outcome any} for this
investigation. An incomplete calculation retains its measured horizon
history and explicitly records the evolution limit.

The \code{ef_evolution_status} record also gives the last accepted and
failed trial timesteps, the smallest active radius, the first cell center
and outer face, and the active rest weight inside these resolution scales.
The raw inner-face values of \(m\), \(mb\), and \(b\) expose the residual
left by the inward constraint solve. An empty central vacuum has \(m=0\)
while the center is regular. A growing residual there must be tested under
radial and temporal refinement. A fixed outer ADM mass alone does not
establish local accuracy.

For an unexcised regular origin, the summary reports the dimensionless
residual \( |m(0)|/M_{\rm ADM} \). When this residual reaches unity,
\path{spatial_resolution_assessment} records
\path{spatial_resolution_lost}. The inferred central error then equals
the entire ADM mass, although regularity requires \(m(0)=0\).
\codesource{Assessing the regular-origin mass residual}
{\path{cpu/src/drivers/einstein_vlasov_wpic.cpp}}
{Reports an order-unity spatial failure separately from the termination cause}
\begin{lstlisting}[language=C++,basicstyle=\ttfamily\scriptsize]
const long double fraction = std::abs(
    static_cast<long double>(field.metric.face_m.front()))
    / field.metric.adm_mass;
summary.ef_spatial_resolution_assessment = fraction >= 1.0L
    ? "spatial_resolution_lost" : "requires_convergence";
\end{lstlisting}
This assessment preserves the accepted state, the evolution equations,
the termination cause, and the \path{interior_evolution_limit}
outcome. Residuals below unity retain the label
\path{requires_convergence}; quantitative accuracy is determined by the
radial and timestep comparisons appropriate to the physical observable.
The same summary gives the smallest active radius divided by the first
radial-cell width. It also reports the radius and contiguous half-maximum
width, in radial cells, of the retained density peak. This density is the
energy density measured by the inward observer defined below. Together,
these quantities show when the contracting matter occupies too few cells
for its density profile to be resolved. The regular-origin assessment
applies when excision is disabled and the ADM mass is positive.

Restart version 2 includes the excision mode in its binary payload and
integrity fingerprint. The interior-controls file also records the mode.
A resume uses the same setting. Version-1 restarts describe enabled excision.
A comparison between modes starts from the same solved null initial slice,
which supplies all trajectories before any excision.


\subsubsection{Density measured by an observer crossing the horizon}
\label{sec:ef-interior-density}

A density is an energy measurement made by a specified observer. The
static observer has density \(T_{vv}/(ab^2)\) in the exterior \(a>0\).
For interior output VENUS instead uses the future radial observer
\[
 u^v=\frac1b,\qquad u^r=\frac{a-1}{2},\qquad
 u^\theta=u^\phi=0.                                    \tag{IEF-I2}
\]
Substitution in the line element gives
\(g_{ab}u^a u^b=-a+(a-1)=-1\), including at \(a=0\) and in \(a<0\).
This defines a local unit timelike observer; a freely falling trajectory
requires a separate geodesic calculation. Its measured particle energy is
\[
 E_u=-p_a u^a=\frac12\left(P+\frac\chi P\right)>0.       \tag{IEF-I3}
\]
The equality follows by substituting (IEF12) and (IEF-I2). Both terms are
positive, which keeps the numerical evaluation well behaved at the horizon.
If marker \(p\) has invariant rest weight \(W_p\), deposits fraction
\(s_{ip}\) in cell \(i\), and that cell has volume \(\Delta V_i\), the output is
\[
 \rho_{{\rm infall},i}
 =\frac1{\Delta V_i}\sum_{p\ {\rm active}}
           s_{ip}W_p\frac{E_{u,p}^2}{P_p}.              \tag{IEF-I4}
\]
This follows from the same null-slice stress measure as (IEF18)--(IEF23),
contracted with \(u^a u^b\). Particle energies are evaluated at marker
positions before the shape weights deposit them into cells.

\codesource{Horizon-crossing observer energy}{\path{cpu/include/ev/weighted_einstein_vlasov_pic_eddington_finkelstein.hpp}}{Positive particle-energy expression used by the interior density deposition}
\begin{lstlisting}[language=C++]
const double energy = 0.5 * (p_r + (1.0 + 2.0 * eta / (radius * radius)) / p_r);
\end{lstlisting}

The profile columns \code{rho_infall} and \code{excision_mask} identify this
observable and its retained region. A center at or inside \(r_{\rm ex}\)
has \code{excision_mask=1}; a particle cloud may still overlap that cell.
Physical plots use \code{excision_mask=0}. History records
\code{active_marker_count}, \code{active_trapped_rest_weight}, and
\code{active_min_radius}. The trapped weight counts active markers with
\(a<0\) at their current positions. A density comparison between coordinate
modules uses the same observer at the same physical events.

\codesource{EF interior visualization}{\path{cpu/tools/render_ef_interior_gif.py}}{Render saved retained-interior density, slope, compactness, and weight budgets}
\begin{lstlisting}[language=bash]
python3 cpu/tools/render_ef_interior_gif.py --run RUN_DIRECTORY \
  --output output/ef_interior.gif --frames 72 --fps 10
\end{lstlisting}

The file \path{ef_interior_controls.txt} records the threshold, buffer, and
fraction used by each fresh run. A restart requires the identical record:
previously removed trajectories determine which interior can still be
evolved. CSV headers are checked before appending. The accepted-step audit
is a separate evidence file with the step and time of the checkpoint.
Its completeness determines which later certification checks are available.

Finite component validation checks observer normalization and stress
contraction, horizon crossing, retained interior trajectories, the trapped
buffer, untrapped-cavity rejection and fixed mass after excision. Synthetic
infall into a prescribed central black hole and collapse from a regular
shell test their respective initial and boundary conditions.

\subsubsection{Bound static energy-cutoff shell initial data}
\label{sec:ef-static-energy-cutoff-shell}

The threshold family starts from a neutral equilibrium with finite radial
support.  Its continuum solver is
\path{cpu/include/ev/weighted_einstein_vlasov_static_shell.hpp}, in namespace
\code{ev::static_einstein_vlasov}.  The distribution uses a single angular momentum magnitude and a cutoff
exponent \(1\le k\le8\). The default \(k=1\) is the neutral shell in
East's static-data construction; \(k>1\) makes the distribution approach
its support edge more smoothly:
\begin{align}
 E(r,w,L)&=e^{\mu(r)}
   \sqrt{1+w^2+\frac{L^2}{r^2}},                    \tag{IEF-S1}\\
 f(E,L)&=C\left(1-\frac{E}{E_0}\right)_+^k
                  \delta(L-L_0),
 \qquad C>0,\quad L_0>0,\quad0<E_0<1.               \tag{IEF-S2}
\end{align}
Here \(L_0\) is the fixed angular momentum magnitude in the shell, \(C\) is the positive
distribution amplitude, \(E_0\) is the cutoff energy, and \(\delta\) is the
Dirac delta distribution.  The subscript \(+\) means
\(x_+=\max(x,0)\).  The energy here is the
dimensionless particle Killing energy per unit rest mass in the static
asymptotically normalized polar metric.  It is not the Newtonian diagnostic
energy and is distinct from the frozen-metric proxy in (IEF58c).
In ordinary language, every particle in the continuum shell has the same
angular momentum magnitude \(L_0\). The factor \((1-E/E_0)^k\)
sets the fall toward zero at the energy cutoff \(E_0\), and the distribution
has support at \(E<E_0\). The exponent \(k\) stays fixed during the
continuum mass rescaling and each comparison of the same solution in both coordinates.  The delta function in
(IEF-S2) expresses the single value \(L=L_0\), rather than a finite
angular momentum width.

Set
\begin{align}
 y(r)&:=\log E_0-\mu(r),& z(r)&:=e^{y(r)},              \tag{IEF-S3a}\\
 W(r)^2&:=z(r)^2-1-\frac{L_0^2}{r^2},&
 \epsilon(r,w)&:=\sqrt{1+w^2+\frac{L_0^2}{r^2}}.   \tag{IEF-S3b}
\end{align}
Matter is present exactly where \(W^2>0\), and then
\[
 |w|<W(r),\qquad
 \Phi(r,w):=1-\frac{\epsilon}{z}
 =\frac{W^2-w^2}{z(z+\epsilon)}>0.                    \tag{IEF-S4}
\]
The second expression for \(\Phi\) is algebraically identical to the first
and avoids subtractive cancellation at the cutoff.

\paragraph{Delta-function Jacobian and continuum sources.}
The Schwarzschild equations integrate over \(\eta=L^2/2\).
Since \(\dd\eta=L\,\dd L\), a shell with \(L_0>0\) satisfies
\[
 \delta(L-L_0)
 =L_0\,\delta(\eta-\eta_0),\qquad \eta_0=\frac{L_0^2}{2}.                  \tag{IEF-S5}
\]
The factor \(L_0\) is essential; omitting it changes every source and the
normalization.  Define
\(p_{t,\Sigma}:=p_\theta+p_\phi=2p_t\) as the sum of the two tangential
pressures and \(n_0\) as the local orthonormal rest-number density.
Substituting (IEF-S2)--(IEF-S5) into the spherical Vlasov moments gives the
complete nonzero orthonormal sources
\begin{align}
 \rho(r)&=\frac{2\pi CL_0}{r^2}
    \int_{-W}^{W}\Phi^k\,\epsilon\,\dd w,                \tag{IEF-S6}\\
 p_{\rm rad}(r)&=\frac{2\pi CL_0}{r^2}
    \int_{-W}^{W}\Phi^k\,\frac{w^2}{\epsilon}\,\dd w,   \tag{IEF-S7}\\
 p_{t,\Sigma}(r)&=\frac{2\pi CL_0}{r^2}
    \int_{-W}^{W}\Phi^k\,
       \frac{L_0^2}{r^2\epsilon}\,\dd w,            \tag{IEF-S8}\\
 n_0(r)&=\frac{2\pi CL_0}{r^2}
    \int_{-W}^{W}\Phi^k\,\dd w,                          \tag{IEF-S9}\\
 j(r)&=\frac{2\pi CL_0}{r^2}
    \int_{-W}^{W}\Phi^k\,w\,\dd w=0.                    \tag{IEF-S10}
\end{align}
Equations (IEF-S6)--(IEF-S10) include both
signs of \(w\); the current vanishes because the unperturbed ansatz is even.
The implementation evaluates these one-dimensional integrals with a fixed
16-point Gauss--Legendre rule on \([-W,W]\).

\paragraph{Static field equations, support edges, and lapse normalization.}
In equilibrium, the enclosed mass grows outward with local energy density,
the lapse gradient supplies the inward relativistic gravitational force, and
the rest-number equation counts invariant particle weight.  The support ends
where the allowed radial-momentum interval shrinks to zero.  With
\(a=e^{-2\lambda}=1-2m/r\), the coupled continuum system is
\begin{align}
 \frac{\dd m}{\dd r}&=4\pi r^2\rho,                    \tag{IEF-S11}\\
 \frac{\dd y}{\dd r}&=-\frac{m/r^2+4\pi r p_{\rm rad}}
                              {1-2m/r}=-\mu_r,          \tag{IEF-S12}\\
 \frac{\dd M_0}{\dd r}&=4\pi r^2 e^\lambda n_0,       \tag{IEF-S13}\\
 \lambda&=-\frac12\log\left(1-\frac{2m}{r}\right).
                                                               \tag{IEF-S14}
\end{align}
The input \code{central_y}, denoted \(y_c>0\), is constant through the exact
inner vacuum.  The angular barrier places its first support edge at
\[
 R_{\rm in}=\frac{L_0}{\sqrt{e^{2y_c}-1}},
 \qquad m(R_{\rm in})=M_0(R_{\rm in})=0,
 \qquad y(R_{\rm in})=y_c.                            \tag{IEF-S15}
\]
The ODE is integrated outward with adaptive step-doubled RK4 until the next
root \(W(R_{\rm out})=0\).  Let \(M=m(R_{\rm out})\) and
\(M_0=M_0(R_{\rm out})\).  Matching to an asymptotically normalized
Schwarzschild exterior fixes the otherwise arbitrary lapse offset and the
cutoff energy:
\[
 \boxed{\ \log E_0=y(R_{\rm out})
       +\frac12\log\left(1-\frac{2M}{R_{\rm out}}\right)\ }.
                                                               \tag{IEF-S16}
\]
Thus \(E_0\) is an output of the normalized solution, not an independently
imposed number.  A finite bound branch must have
\(R_{\rm out}>2M\), \(0<E_0<1\), and
\[
 E_{\rm bind}=M_0-M>0,\qquad
 f_{\rm bind}=\frac{M_0-M}{M_0}>0.                    \tag{IEF-S17}
\]
In the exact inner vacuum,
\(m=M_0=0,\lambda=0,\mu=\log E_0-y_c\).  In the exact exterior,
\begin{align}
 m&=M,&M_0&=\hbox{constant},&
 e^{2\mu}&=1-\frac{2M}{r},&e^{2\lambda}&=
       \left(1-\frac{2M}{r}\right)^{-1}.              \tag{IEF-S18}
\end{align}
The continuum solution also checks (IEF58i) directly.  Its reported
constraint defect and relative Jeans residual are numerical errors, not
additional equations used to close the system.

\paragraph{Exact scaling family.}
If \(s>0\), a solved equilibrium can be rescaled without re-shooting:
\begin{align}
 (r,m,M,M_0,L_0,R_{\rm in},R_{\rm out})
 &\longmapsto
 s(r,m,M,M_0,L_0,R_{\rm in},R_{\rm out}),           \tag{IEF-S19}\\
 C&\longmapsto C/s,\notag\\
 (\rho,p_{\rm rad},p_{t,\Sigma},n_0)&\longmapsto
 s^{-2}(\rho,p_{\rm rad},p_{t,\Sigma},n_0).           \tag{IEF-S20}
\end{align}
The dimensionless functions \(y,\mu,\lambda\), \(E_0\), compactness, and
\(f_{\rm bind}\) are unchanged.  The public
\code{rescale_to_adm_mass} routine uses \(s=M_{\rm target}/M\).  This
continuum rescaling selects the physical length/mass scale of the baseline;
the later marker-level fixed-ADM normalization in (IEF59a) separately removes
finite-sampling and perturbed-family mass changes.

\paragraph{Invariant WPIC weights and quiet sampling.}
The reduced invariant phase measure can be written
\[
 \dd\Gamma_{\rm inv}
 =8\pi^2e^\lambda\,\dd r\,\dd w\,\dd\eta.
                                                               \tag{IEF-S21}
\]
After applying (IEF-S5), a positive-\(w\) quadrature point representing
\(\Delta r\,\Delta w\) receives
\[
 W_p=8\pi^2 CL_0 e^{\lambda(r_p)}
       \Phi(r_p,w_p)^k\,\Delta r\,\Delta w,              \tag{IEF-S22}
\]
and an adjacent \(-w_p\) marker receives exactly the same weight, radius, and
\(\eta=L_0^2/2\).  These are already the invariant weights of (EWP3):
passing them through
\path{canonicalize_weighted_einstein_vlasov_pic_initial_weights}
again would multiply by \(e^\lambda\) twice and is forbidden.

The library supplies two sampling rules.  Tensor Gauss--Legendre quadrature
is used for deterministic convergence tests.  The driver uses a dedicated
shifted rank-one rule with exactly
\(N_{\rm marker}=2N_{\rm pair}\) markers.  For deterministic
Cranley--Patterson shifts \((\Delta_r,\Delta_w)\), a generator \(g\) coprime
to \(N_{\rm pair}\), and fractional part \(\{\cdot\}\), it samples
\begin{align}
 u_{r,k}&=\left\{\frac{k+1/2}{N_{\rm pair}}+\Delta_r\right\},&
 u_{w,k}&=\left\{\frac{(gk\bmod N_{\rm pair})+1/2}
                         {N_{\rm pair}}+\Delta_w\right\},       \tag{IEF-S23}\\
 r_k&=R_{\rm in}+(R_{\rm out}-R_{\rm in})u_{r,k},&
 w_k&=W(r_k)u_{w,k},                                    \tag{IEF-S24}\\
 \Delta r\,\Delta w&=
 \frac{(R_{\rm out}-R_{\rm in})W(r_k)}{N_{\rm pair}}.
                                                               \tag{IEF-S25}
\end{align}
Each \(k\) emits adjacent \((+w_k,-w_k)\) records.  Their persistent
\code{id} and \code{source_cell} are unique and equal record by record, so
the coordinate-map manifest can authenticate the empirical measure without
confusing a source identity with a radial deposition bin.

Changing \code{--high-sampling-realization} changes both shifts and therefore
the locations and unequal quadrature weights, while retaining exact pairwise
current cancellation before a kick.  Lattices with different
\(N_{\rm pair}\) are not nested, and two shifts are only the implemented
minimum PIC-discreteness cross-check; they are not a calibrated statistical
confidence interval.  Tensor-rule quadrature orders do not control a
fixed-size shifted lattice and are consequently not exposed by the
production driver.  The full continuum support must also leave at least four
vacuum mesh cells at each radial boundary so the TSC stencil never clips the
shell.

The essential library calls are:
\codesource{Static Einstein--Vlasov shell construction}{\path{cpu/include/ev/weighted_einstein_vlasov_static_shell.hpp}}{Continuum solve, exact mass rescaling, and shifted-stratified marker sampling}
\begin{lstlisting}[language=C++,basicstyle=\ttfamily\scriptsize]
namespace static_ev = ev::static_einstein_vlasov;

static_ev::Parameters p;
p.amplitude = C;
p.angular_momentum = ell0;
p.central_y = y_c;
p.cutoff_power = k;  // 1 <= k <= 8; choose 2 for a smoother support edge.
auto continuum = static_ev::solve(p);
continuum = static_ev::rescale_to_adm_mass(continuum, target_M);

static_ev::MarkerSamplingOptions sampling;
sampling.rule = static_ev::MarkerSamplingRule::shifted_stratified;
sampling.shifted_pair_count = marker_count / 2;
sampling.realization = realization;
auto sampled = static_ev::sample_markers(continuum, sampling);
// sampled.markers already carry invariant Schwarzschild WPIC weights.
\end{lstlisting}

\paragraph{Controlled inward perturbation.}
The exact equilibrium baseline has \(s_L=1\) and zero inward kick.  To ask a
collapse-versus-halo question without starting from an unbound Gaussian
cloud, the production family applies one fixed nonnegative kick after the
continuum quadrature:
\[
 w_p\longmapsto w_p-\kappa\,{\cal K}(r_p),\qquad
 {\cal K}(r)=
 \begin{cases}
 1,&r_c=0,\\[2pt]
 r^2/(r^2+r_c^2),&r_c>0.
 \end{cases}                                           \tag{IEF-S26}
\]
The controls are \code{--high-radial-infall} \(=\kappa\) and
\code{--high-infall-core-radius} \(=r_c\).  Each adjacent pair then has exact
mean radial momentum \(-\kappa{\cal K}(r)\); weights, radii, \(\eta\), and IDs
are unchanged.  A nonzero kick, or a subsequent \(s_L\ne1\) rescaling, is
deliberately nonequilibrium and must not be described as the original static
solution.  The same fixed \((\kappa,r_c)\) is used at every \(s_L\), and the
common weight factor in (IEF59a) is applied only after both perturbations.

The corresponding driver interface is
\codesource{Static-shell production interface}{\path{cpu/src/drivers/einstein_vlasov_wpic.cpp}}{High-compactness initial-condition, kick, angular-scale, and target-mass options}
\begin{lstlisting}[basicstyle=\ttfamily\scriptsize]
--high-initial-condition energy_cutoff_shell
--high-energy-cutoff-amplitude C
--high-energy-cutoff-ell0 ELL0
--high-energy-cutoff-central-y Y_C
--high-energy-cutoff-power K
--high-markers N_EVEN
--high-sampling-realization REALIZATION
--high-radial-infall KAPPA
--high-infall-core-radius R_C
--high-angular-momentum-scale S_L
--high-target-initial-adm-mass M_TARGET
\end{lstlisting}
For this IC, \code{--high-markers} must be even and
\code{--high-specific-angular-momentum} is rejected: the baseline
\(L_0\) comes from \code{--high-energy-cutoff-ell0}, and only
\code{--high-angular-momentum-scale} defines the scan.  A saved physical IC
is regenerated from its continuum, sampler, kick, and realization controls
and fingerprint-compared before reuse, including on resume.

The JSON summary keeps the configured and rescaled \(C,L_0,y_c\), unscaled
and continuum ADM masses, \(E_0,R_{\rm in},R_{\rm out},M_0\), binding energy
and fraction, central lapse, maximum compactness, constraint/Jeans residuals,
pair count, represented rest weight, rest-weight quadrature error, and maximum
\(E/E_0\).  Component validation checks source normalization, finite support and both
vacuum regions, scale covariance, constraint and Jeans residuals, marker
support and IDs, exact zero-current pairs, sampling/deposition refinement,
independent shifted realizations and a short self-consistent equilibrium
evolution. Long-time nonlinear stability requires evolution over several
dynamical times together with resolution, realization, timestep and domain
comparisons.

\subsubsection{Diagnostics and scientific acceptance checks}

Every production run should record at least:
\begin{itemize}
\item \(M_{\rm ADM}\), total represented rest weight, active, excised, and
      escaped rest weights, and deposition roundoff;
\item \(\min a\), \(\max(2m/r)\), the radius of that maximum, and the
      outermost apparent-horizon radius;
\item \(\min b\) and \(\min P\), which must remain strictly positive;
\item the first advanced time at which \(a=0\), post-horizon evolution time,
      and excision radius;
\item Lagrangian radii, density and mass profiles versus areal radius, and
      the mass-weighted mean, root-mean-square (RMS), and maximum specific
      angular momentum;
\item independent residuals of (IEF28)--(IEF31) when a dedicated validation
      or post-processing artifact supplies the needed derivatives; these are
      not ordinary production-history columns;
\item the authenticated null-map analytic-extension error for (IEF54); the
      ordinary evolved profile does not currently emit a generic exterior
      vacuum-error column;
\item thread-count reproducibility and convergence with radial cells, marker
      count, particle shape, timestep, and characteristic substeps.
\end{itemize}
No arbitrary floor should replace \(a\), \(b\), or \(P\) in the equations.
Loss of \(b>0\), \(P>0\), finite fields, or fixed-\(M\) bookkeeping is a failed
evolution, not a physical black hole signal.

\paragraph{Fail-closed black hole candidate and certificate.}
An apparent-horizon flag by itself is deliberately not the final
black hole classification.  Let \(v_H\) be the first saved advanced time at
which the outer apparent horizon is found, let \(v_f\) be the last saved
advanced time, and let \(M_f=M_{\rm family}\) be the fixed-family ADM mass.
Define
\begin{align}
 \Delta v_{\rm post}&=v_f-v_H,&
 N_{>H}&=\#\{i:v_i>v_H\},&
 N_H&=\#\{i:v_i\ge v_H\ \hbox{and}\ r_H(v_i)>0\}.
                                                               \tag{IEF57c}
\end{align}
The symbol \(\#\) counts the saved rows satisfying the condition in braces.
The configured floors cannot be weakened below their defaults.  The intrinsic
post-horizon duration and sampling gates are
\[
 \boxed{\ \Delta v_{\rm post}/M_f\ge5,\qquad N_{>H}\ge3,
         \qquad N_H\ge1\ }.                            \tag{IEF57d}
\]
The independently stored value of \(\Delta v_{\rm post}\) must agree with
\(v_f-v_H\) to the configured end-time tolerance.  The run must use ingoing
EF coordinates, report a finite \(v_H\ge0\), set both
\path{horizon_formed} and \path{evolved_past_horizon}, report the raw driver
outcome \path{black_hole}, and reach
\(\max_i{\cal C}_{\max}(v_i)\ge0.99\) by default.  These conditions prevent
a single marginal output row from being called a resolved continuation.

Output cadence alone is too sparse to show that the characteristic equations
remained regular between saved rows.  The driver therefore audits the initial
state and every accepted EF step.  If \({\cal A}\) is that set of audited
states, it records
\begin{align}
 P_{\min}^{\rm acc}
  &=\min_{k\in{\cal A}}\min_p P_{p,k}>0,&
 b_{\min}^{\rm acc}
  &=\min_{k\in{\cal A}}\min_{i,\,x\in\{\mathrm{center},\mathrm{face}\}}
       b_{i,x,k}>0.                                    \tag{IEF57e}
\end{align}
Every marker value and every stored metric, metric-derivative, mass,
compactness, and source array must be finite at those states.  At each audit
the radial constraint must also have used the fixed EF ADM anchor
\(M_{\rm EF}=M_{\rm null}\) imported with the null slice and satisfy
\[
|M_{\rm metric,k}-M_{\rm EF}|
 \le4096\,\epsilon_{\rm mach}\max(1,|M_{\rm EF}|).    \tag{IEF57f}
\]
Here \(M_{\rm metric,k}\) is the ADM mass stored by the metric solve at
audited state \(k\), and \(\epsilon_{\rm mach}\) is the unit roundoff of the
double-precision floating-point type.
The complete extrema and Boolean checks are persisted atomically in
\path{ef_post_horizon_audit.state}, together with the exact accepted step and
advanced time.  Sidecar version 2 includes the cumulative local metric check
\code{causal_excision_valid_throughout}.  A version-1 file can restore its
other extrema; its weaker excision record marks the accumulated audit
incomplete.  A resume accepts the sidecar only when those two values match
the restart.  A missing, stale, or old-format sidecar marks the audit
incomplete and makes black hole promotion impossible.  Despite the sidecar's
name, these extrema cover the whole EF evolution, not merely the saved
post-horizon rows; this is what establishes that positivity and finite fields
were maintained \emph{through} horizon formation.
The summary must expose schema
\path{venus_wpic_ef_post_horizon_evidence_v1}, set
\path{available=true} and \path{accepted_step_audit_complete=true}, and report
a strictly positive \path{accepted_step_audit_samples}; absent fields fail
closed.

The accepted-step call is deliberately inside the advance loop, before the
optional output-cadence branch:
\codesource{EF accepted-step audit}{\path{cpu/src/drivers/einstein_vlasov_wpic.cpp}}{Per-accepted-state positivity, finiteness, mass-anchor, and excision evidence}
\begin{lstlisting}[language=C++,basicstyle=\ttfamily\scriptsize]
field = result.endpoint_field;
advanced_time += result.accepted_dt;
++step;
update_ef_post_horizon_audit(
    post_horizon_audit, markers, field, conserved_adm_mass,
    true, evolution.excision_a_threshold,
    evolution.excision_buffer_cells);
// Restart and audit sidecar are written together at checkpoints.
\end{lstlisting}
Thus increasing \path{high_output_steps} cannot make a failed intermediate
state invisible to the candidate gate.

The analyzer separately verifies four positive mass anchors:
\(M_{\rm target}\), the requested family mass; \(M_f\), the realized fixed
family mass; \(M_{\rm null}\), the imported null-slice mass; and
\(M_{\rm dep}\), the mass reconstructed from the initial EF deposition.  It
then forms
\begin{align}
 \delta_T&=\frac{|M_f-M_{\rm target}|}{M_{\rm target}},&
 \delta_D&=\frac{|M_{\rm null}-M_f|}{M_f},&
 \delta_{\rm map}&=\frac{|M_{\rm dep}-M_{\rm null}|}{M_{\rm null}},
                                                               \tag{IEF57g}\\
 \delta_H^M&=\max_i\frac{|M_{\rm ADM}(v_i)-M_f|}{M_f},&
 \delta_H^W&=\max_i
 \frac{|W_{\rm tot}-W_{\rm active}-W_{\rm exc}-W_{\rm esc}|_i}
      {\max(|W_{\rm tot,i}|,10^{-300})}.              \tag{IEF57h}
\end{align}
In (IEF57h), \(W_{\rm tot}\), \(W_{\rm active}\), \(W_{\rm exc}\), and
\(W_{\rm esc}\) are respectively the total, active, excised, and escaped
invariant rest weights at a saved state.
The default bounds are
\(\delta_T\le10^{-8}\),
\(\delta_D,\delta_{\rm map},\delta_H^M\le10^{-2}\), and
\(\delta_H^W\le10^{-12}\), in addition to the analyzer's ordinary strict
summary-quality checks.  Thus excised rest weight is accounted for rather
than discarded, while the ADM anchor is not redefined after excision.
Those ordinary summary checks are also explicit:
\[
 \frac{|M_{\rm final}-M_{\rm initial}|}{M_{\rm initial}}<10^{-2},
 \qquad
 \frac{|W_{\rm final}-W_{\rm initial}|}{W_{\rm initial}}<10^{-12}.
                                                               \tag{IEF57h-1}
\]
The strict inequalities here match the general run-quality gate; they are in
addition to, not substitutes for, the all-history bounds in (IEF57h).  A
mapped EF candidate must also pass the authenticated null-map preflight and
the fixed-ADM/IC provenance checks described above.

The excision controls must be explicit, with
\(a_{\rm ex}<0\) and an integer buffer \(n_{\rm buf}\ge1\).  Every saved row
with a positive excision radius must obey
\[
 r_{\rm ex}(v_i)>0\quad\Longrightarrow\quad
 v_i>v_H\ \hbox{and}\ 0<r_{\rm ex}(v_i)<r_H(v_i),    \tag{IEF57i}
\]
and at least one such causal row must have \(W_{\rm exc}>0\).  The
accepted-step audit must independently record that a causal excision was
seen, that positive interior rest weight was actually excised, and that
\code{causal_excision_valid_throughout} remains true.  At every audited
state with excision, this last condition evaluates (IEF58), positive finite
\(b\), and the entire outward buffer in the reconstructed metric.  A nominal
excision setting that never removes matter therefore supplies no
post-horizon evidence.

Passing all the preceding tests gives only the machine classification
\path{black_hole_candidate}.  It is independent of the positive-binding
condition used for a halo: binding is part of the chosen campaign family, not
a logical prerequisite for recognizing a trapped region.  A candidate is
promoted to \path{black_hole_certified} only by a separate, SHA-256
authenticated
\path{venus_wpic_excision_sensitivity_certificate_v1}.  The analyzer hashes
and re-reads every referenced summary instead of trusting the certificate's
claimed conclusion.  At least two distinct repeats are required; every repeat
must itself pass the complete intrinsic candidate gate and retain the same
physical-IC hash, base and realized IC fingerprints, coordinate-map-manifest
hash, angular momentum, and fixed-family ADM mass.  The primary run and every
repeat must also carry independently verified \path{binary.sha256} and
\path{source.sha256} records.  The certificate names the expected executable
and source-manifest hashes, and the analyzer recomputes them from disk and
requires exact agreement.  In plain language, an excision-sensitivity repeat
cannot certify a result obtained with a different compiled program or source
tree.  At least one repeat must
differ in \(a_{\rm ex}\) while holding \(n_{\rm buf}\) fixed, and at least
one must differ in \(n_{\rm buf}\) while holding \(a_{\rm ex}\) fixed.  The
scientific repeat protocol holds every other run control fixed; the analyzer
directly authenticates the physical, map, executable, and source provenance
listed above.

For the primary run plus all authenticated repeats, the sensitivity agreement
conditions are
\begin{align}
 \frac{\max v_H-\min v_H}{M_f}&\le1,&
 \frac{\max r_H-\min r_H}{\operatorname{median}(r_H)}&\le0.05,
                                                               \tag{IEF57j}
\end{align}
Here each \(v_H\) is the run's first trapped time, whereas each \(r_H\) is the
positive \path{apparent_horizon_radius} reported by that run's final summary;
the analyzer does not silently substitute a radius from another output row.
Command-line choices may make these bounds tighter but not looser.  Missing
authentication, a missing independent
variation, changed physical provenance, a repeat that is not itself a
candidate, or failure of either spread test leaves the primary run a candidate
only.  The final angular momentum bracket accepts
\path{black_hole_certified}, never a bare candidate; a candidate plus a halo
is reported as
\path{inconclusive_due_to_uncertified_black_hole_candidate}.

The analyzer keeps that promotion explicit:
\codesource{Black-hole certification analyzer}{\path{cpu/tools/analyze_einstein_vlasov_wpic.py}}{Candidate classification and authenticated excision-sensitivity promotion}
\begin{lstlisting}[language=Python,basicstyle=\ttfamily\scriptsize]
classification = "black_hole_candidate"  # intrinsic gate passed
sensitivity = excision_sensitivity_gate(summary_path, record, args)
if sensitivity["certifies_black_hole"]:
    classification = "black_hole_certified"
\end{lstlisting}

The frequently used ``Lagrangian radii'' are rest-weight quantiles, not new
coordinates.  At a saved time \(t\), sort the active markers by radius and
define \(r_q(t)\) as the smallest marker radius satisfying
\[
 \sum_{p\,{\rm active}:\,r_p\le r_q}W_p
 \ge q\sum_{p\,{\rm active}}W_p,
 \qquad q\in\{0.1,0.5,0.9\}.                           \tag{IEF57a}
\]
Thus \(r_{50}\) encloses half of the active rest weight and \(r_{90}\)
encloses 90 percent.  The saved maximum compactness is
\[
 {\cal C}_{\max}(t):=\max_r\frac{2m(t,r)}r.           \tag{IEF57b}
\]
These definitions are identical in the SC and EF drivers wherever the two
descriptions overlap.

\paragraph{Initially bound family diagnostic and mass schema.}
The angular scan is defined on the normalized Schwarzschild \(\tau=0\)
family, not on the later extracted null slice.  Let \(W_p\) be the invariant
rest weight of marker \(p\).  VENUS records
\begin{align}
 M_0&=\sum_p W_p,
 &E_{\rm bind}&=M_0-M_{\rm ADM}^{\rm family},             \tag{IEF58a}\\
 f_{\rm bind}&=\frac{E_{\rm bind}}{M_0}
 =\frac{M_0-M_{\rm ADM}^{\rm family}}{M_0}.              \tag{IEF58b}
\end{align}
Here \(M_0\) is
\code{fixed_family_rest_weight} and \(f_{\rm bind}\) is
\code{fixed_family_binding_fraction}.  Both total rest weight and ADM mass
are invariant under the asymptotically normalized SC--EF coordinate change,
so the SC and mapped-EF summaries must contain the same values.  The driver
stores the numbers in the comma-separated-value (CSV) file
\path{fixed_family_initial_bound_diagnostics.csv}; a resume reads and
algebraically validates that immutable one-row sidecar instead of inferring
the family from the evolved checkpoint.  The condition \(f_{\rm bind}>0\) is
necessary global positive-binding-energy evidence for a halo label, but is
not sufficient for virialization: a globally bound system may undergo broad
breathing oscillations, lose a component, or remain unresolved on the saved
time interval.  Accordingly the summary boolean is conservatively named
\code{fixed_family_positive_binding_energy}, not ``all particles bound.''

For additional transparency the code evaluates, on this same SC \(\tau=0\)
slice,
\begin{align}
 \epsilon_p&=\sqrt{1+w_p^2+2\eta_p/r_p^2},
 &E_{p,\mathrm{proxy}}&=e^{\mu(0,r_p)}\epsilon_p,          \tag{IEF58c}\\
 \langle E_{\mathrm{proxy}}\rangle_{W}
 &=\frac{1}{M_0}\sum_pW_pE_{p,\mathrm{proxy}},
 &F_{E\ge1}&=\frac{1}{M_0}\sum_pW_p
       \mathbf 1[E_{p,\mathrm{proxy}}\ge1].              \tag{IEF58d}
\end{align}
The indicator \(\mathbf 1[\mathcal P]\) equals one when proposition
\(\mathcal P\) is true and zero otherwise.
It reports \(\min_pE_{p,\mathrm{proxy}}\),
\(\max_pE_{p,\mathrm{proxy}}\), the weighted mean in (IEF58d), and
\(F_{E\ge1}\).  In a static Schwarzschild spacetime
\(-u_\tau=e^\mu\epsilon\) is the Killing energy per unit rest mass.  The
collapse data are dynamical and have no such timelike Killing symmetry, so
(IEF58c) is only an instantaneous frozen-metric energy-at-infinity proxy; it
is not conserved and \(F_{E\ge1}\) is reported but never used as a hard halo
cut.  Likewise the history column called \code{virial_ratio} is the Newtonian
proxy \(2K/|W|\), not a relativistic virial identity.  The sustained radius,
metric, general-relativistic (GR) stress-balance, escape, and time-series
tests therefore remain
essential, and an unresolved breathing state remains \code{undetermined}.

Three distinct ADM quantities appear in every mapped run and must not be
collapsed into one field:
\begin{align}
 M_{\rm family}&=M_{\rm ADM}(\tau=0),                    \tag{IEF58e}\\
 M_{\rm null}&=M_{\rm ADM}\text{ at the extracted null slice},\tag{IEF58f}\\
 M_{\rm dep}&=M_{\rm ADM}\text{ reconstructed from EF markers at }\widehat v=0.
                                                               \tag{IEF58g}
\end{align}
Their exact schema mapping is \(M_{\rm family}\leftrightarrow\)
\path{fixed_family_adm_mass}, \(M_{\rm null}\leftrightarrow\)
\path{null_slice_conserved_adm_mass}, and \(M_{\rm dep}\leftrightarrow\)
\path{initial_deposited_adm_mass}.
The family-to-null donor conservation error and null-to-deposition map error
are reported separately as
\begin{equation*}
 \delta_{\rm donor}=\frac{|M_{\rm null}-M_{\rm family}|}{M_{\rm family}},
 \qquad
 \delta_{\rm map}=\frac{|M_{\rm dep}-M_{\rm null}|}{M_{\rm null}}.\tag{IEF58h}
\end{equation*}
Their schema names are
\path{null_slice_adm_relative_drift_from_family} and
\path{initial_deposition_adm_relative_error_from_null_slice}.
Neither numerical mismatch redefines the fixed-ADM family.

\paragraph{Relativistic late-time balance diagnostic.}
A stationary, spherically symmetric anisotropic distribution in Schwarzschild
variables obeys the radial Einstein--Vlasov anisotropic
Tolman--Oppenheimer--Volkoff (TOV), or Jeans, balance
equation
\begin{equation*}
 {\cal R}_{J}(r)
 =\frac{d p_{\rm rad}}{dr}+(\rho+p_{\rm rad})\mu_r
  -\frac{p_{t,\Sigma}-2p_{\rm rad}}{r}=0,
 \qquad p_{t,\Sigma}=p_\theta+p_\phi=2p_t .              \tag{IEF58i}
\end{equation*}
The three terms are the radial pressure gradient, gravitational force
density, and force density from pressure anisotropy.  Their cancellation in
(IEF58i) gives the instantaneous GR static-balance test.  The Newtonian
\code{virial_ratio} measures a different approximation.  A small balance
residual over a sustained interval is necessary for stationarity, while
nonlinear stability also requires following perturbations over the relevant
dynamical times.

For cell centers \(r_i\), VENUS uses the centered nonuniform-grid derivative
\((p_{{\rm rad},i+1}-p_{{\rm rad},i-1})/(r_{i+1}-r_{i-1})\).
Denote this discrete radial derivative by \(D_r\).  The code records
\begin{align}
 {\cal J}_{1}
 &=\frac{\displaystyle\sum_{i\in\Omega}\Delta V_i
                       |{\cal R}_{J,i}|}
        {\displaystyle\sum_{i\in\Omega}\Delta V_i
          \left(\left|D_rp_{{\rm rad},i}\right|
          +\left|(\rho_i+p_{{\rm rad},i})\mu_{r,i}\right|
          +\left|\frac{p_{t,\Sigma,i}-2p_{{\rm rad},i}}{r_i}\right|\right)},
                                                               \tag{IEF58j}\\
 \Delta V_i&=\frac{4\pi}{3}
       \left(r_{i+1/2}^{3}-r_{i-1/2}^{3}\right).              \tag{IEF58k}
\end{align}
The common volume weight makes \({\cal J}_1\) dimensionless; zero is exact
balance, while a value of order unity means that the uncancelled residual is
comparable to the magnitudes of the force terms.  The support set \(\Omega\)
excludes the first and last derivative cells, any cell whose three-point
stencil is invalid, and cells with
\(\rho_i\le10^{-6}\max_j\rho_j\).  At least three supported cells are
required.  Otherwise the history records \(-1\), explicitly meaning
``unavailable,'' never zero.

SC outputs supply \(\rho,p_{\rm rad},p_{t,\Sigma},\mu_r\) directly.  Before
a horizon,
an EF output reconstructs the same-event Schwarzschild quantities as
\begin{equation*}
 \rho=\frac{T_{vv}}{ab^2},\qquad
 p_{\rm rad}=\rho-S,\qquad
 p_{t,\Sigma}=\frac{2T_{\theta\theta}}{r^2},\qquad
 \mu_r=\frac{m+4\pi r^3p_{\rm rad}}{r^2a},
 \qquad a=1-\frac{2m}{r}.                                  \tag{IEF58l}
\end{equation*}
These are the same local tensor pullbacks used by the coordinate-parity
audit.  The pressure derivative used in an ordinary EF output, however,
follows its constant-\(v\) radial array.  The map in (IEF38c-6) gives
\[
 \left.\partial_rp_{\rm rad}\right|_v
 =\left.\partial_rp_{\rm rad}\right|_\tau
  -h\,\partial_\tau p_{\rm rad},\qquad
 h=e^{\lambda-\mu}.                                    \tag{IEF58l-1}
\]
Thus the native EF residual approaches the Schwarzschild static-balance
residual in a stationary state.  During time-dependent motion it also
contains the null-slice pressure variation, so it is used as a stationarity
proxy together with the radius and metric histories.  An exact dynamical
cross-coordinate residual requires reconstructing a common spacelike slice
or including the time derivative in (IEF58l-1).
Formula (IEF58l) is evaluated
only where \(a>0\) and \(b>0\); once the polar chart fails at a trapped surface,
the EF Jeans value is unavailable even though the EF evolution itself remains
regular and continues through the horizon.

The implementation mirrors the displayed terms without a hidden floor:
\codesource{Relativistic Jeans diagnostic}{\path{cpu/src/drivers/einstein_vlasov_wpic.cpp}}{Profile-gradient, gravity, anisotropy, and normalized residual calculation}
\begin{lstlisting}[basicstyle=\ttfamily\scriptsize]
dp_dr      = (p_r[i+1] - p_r[i-1]) / (r[i+1] - r[i-1]);
gravity    = (rho[i] + p_r[i]) * mu_r[i];
anisotropy = (p_t_sum[i] - 2*p_r[i]) / r[i];
R_J        = dp_dr + gravity - anisotropy;
J1         = sum(dV*abs(R_J))
           / sum(dV*(abs(dp_dr)+abs(gravity)+abs(anisotropy)));
\end{lstlisting}
\paragraph{Exact time-series halo gates.}
The halo classifier tests the decrease in half-rest-mass radius, sustained
radius and metric behavior, and relativistic stress balance.  The following definitions specify its
complete default criteria.  Let
\({\cal S}_{\rm hist}:=\{t_i\}\) be the complete
saved history, let \(t_0,t_f\) be its first and last times, and let
\(T_{\rm cfg}\) be the requested end time.  The late set \({\cal L}\) is
\[
 t_i\in{\cal L}
 \quad\Longleftrightarrow\quad
 t_i\ge t_f-0.30(t_f-t_0),
 \qquad N_{\cal L}\ge5.                                \tag{IEF58m}
\]
The history must contain at least two strictly time-ordered samples; every
\(r_{50},r_{90}\) must be positive and every saved compactness and Jeans
value must be finite.  A Jeans value of \(-1\) remains a finite, explicit
``unavailable'' marker and is excluded only from the Jeans subset below.

For any saved scalar series \(x_i\), write \(Q_p[x]\) for its linearly
interpolated percentile after sorting: its zero-based sorted position is
\(p(N-1)\).  Write
\[
 B_{\cal I}[x]
 =\frac{\sum_{i\in{\cal I}}(t_i-\bar t_{\cal I})
                             (x_i-\bar x_{\cal I})}
        {\sum_{i\in{\cal I}}(t_i-\bar t_{\cal I})^2},
 \qquad
 \Delta t_{\cal I}=\max_{i\in{\cal I}}t_i-
                    \min_{i\in{\cal I}}t_i,           \tag{IEF58n}
\]
where bars are ordinary sample means.  Thus \(B\) is the least-squares
straight-line slope and \(B\Delta t\) is its fitted change across the stated
window.

Before any late-time shape test, all of the following gates must pass:
\begin{align}
 t_f&\ge(1-10^{-6})T_{\rm cfg},&
 f_{\rm bind}&>0,&
 \hbox{no apparent horizon},                           \tag{IEF58o}\\
 \frac{W_{\rm escaped}}{W_{\rm initial}}&\le0.05,&
 \min_{t_i\in{\cal S}_{\rm hist}}r_{50}(t_i)&<0.95r_{50}(t_0),&
 \max_{t_i\in{\cal S}_{\rm hist}}{\cal C}_{\max}(t_i)&<0.99. \tag{IEF58p}
\end{align}
Here \(W_{\rm escaped}\) and \(W_{\rm initial}\) are escaped and initial
invariant rest weights.  The compression test requires the perturbed object
to contract measurably before settling; it prevents an almost unevolved
initial state from being called a newly formed halo.  The last condition
prevents a Schwarzschild run that merely stalls next to its coordinate
singularity from being mislabeled as a halo.

With \(\delta=10^{-300}\) used only to make a formally safe denominator, the
late radius and metric gates are exactly
\begin{align}
 Q_{0.5}[r_{50}]&\le1.50r_{50}(t_0),&
 \frac{Q_{0.9}[r_{50}]-Q_{0.1}[r_{50}]}
      {\max(Q_{0.5}[r_{50}],\delta)}&\le0.75,          \tag{IEF58q}\\
 \left|B_{\cal L}[r_{50}]\frac{\Delta t_{\cal L}}
                                  {r_{50}(t_0)}\right|&\le0.25,&
 \left|B_{\cal L}[r_{90}]\frac{\Delta t_{\cal L}}
                                  {r_{90}(t_0)}\right|&\le0.25,\tag{IEF58r}\\
 Q_{0.9}[{\cal C}_{\max}]-Q_{0.1}[{\cal C}_{\max}]&\le0.05,&
 \left|B_{\cal L}[{\cal C}_{\max}]\Delta t_{\cal L}\right|&\le0.05.
                                                               \tag{IEF58s}
\end{align}
In plain language, the inner half-mass radius must remain bounded, neither
the half-mass nor 90-percent radius may show sustained growth or shrinkage,
and the strong-field geometry must stop drifting.

Let \({\cal L}_J\subseteq{\cal L}\) contain only samples with
\({\cal J}_1\ge0\).  The balance gates are
\begin{align}
 N_{{\cal L}_J}&\ge
 \max\!\left(3,\left\lceil0.75N_{\cal L}\right\rceil\right),&
 Q_{0.5}[{\cal J}_1]&\le0.25,&
 Q_{0.9}[{\cal J}_1]&\le0.25,                         \tag{IEF58t}\\
 B_{{\cal L}_J}[{\cal J}_1]\Delta t_{{\cal L}_J}&\le0.05.&&
                                                               \tag{IEF58u}
\end{align}
The final gate limits an \emph{increase}; it deliberately has no absolute
value, so an improving (negative-slope) balance residual is allowed.  A
horizonless state that fails any gate is ``undetermined'' or ``dispersal'' as
appropriate, never a certified halo.

The SC and EF drivers call the same implementation in
\path{include/ev/weighted_einstein_vlasov_halo_classifier.hpp}; at run
finalization it reads the complete saved history and applies all of the cuts
above.  The complete audit is written under the
\path{high_compactness.halo_classifier} summary object, including every
threshold, late-window statistic, Boolean gate, and stable
\path{failed_checks} name.  Only an empty failed-check list sets
\path{virialized_halo=true}, emits \path{outcome=virialized_halo}, and allows
\code{--expected-outcome halo} to return success.  The Python threshold
analyzer independently recomputes the time-series evidence and verifies that
its result is consistent with this driver audit.  Analyzer options may impose
additional sensitivity cuts, but cannot turn a driver-rejected state into a
halo.  Every numeric cut still requires resolution and static-baseline
sensitivity tests: these numbers are conservative analysis choices, not new
laws of relativistic dynamics.  No band in \(2K/|W|\) is a halo gate.

\subsubsection{Validation hierarchy and Schwarzschild overlap experiment}

Component validation covers continuum identities, momentum maps,
constraints, characteristics, horizons and excision. A scientific calculation
extends these checks through the following hierarchy:
\begin{enumerate}
\item \textbf{Minkowski deposition and characteristics.}
      Verify (IEF47)--(IEF50), \(a=b=1\), exact radial characteristics, and
      conservation of the flat-space Hamiltonian for \(L\ne0\).
\item \textbf{Momentum-map round trip.}
      Check (IEF42)--(IEF44) from strongly inward to strongly outward \(w\),
      including cancellation-prone inward data.
\item \textbf{Vacuum Schwarzschild.}
      With fixed \(M\) and no active markers, recover (IEF54), locate
      \(r_H=2M\), and verify finite characteristic derivatives at and inside
      \(r=2M\).
\item \textbf{Constraint and parallel identities.}
      Verify (IEF30), \(b>0\), fixed outer \(Y=M\), deterministic deposition,
      and identical frozen pushes for different thread counts.
\item \textbf{Horizon and excision bookkeeping.}
      Cross \(a=0\), excise only behind the horizon, retain total rest weight,
      and retain fixed ADM mass after source removal.
\item \textbf{Static cross-coordinate solution.}
      Transform one equilibrium Schwarzschild marker ensemble
      using (IEF40)--(IEF43).  Compare Misner--Sharp mass, stress tensor,
      areal density, and stationary orbits at increasing resolution.
\item \textbf{Dynamical same-event coordinate equivalence.}
      Generate the EF initial state by solving (IEF45m) on the dense SC donor
      history and by intersecting every marker worldline with that solved null
      surface.  Solve the inverse eikonal (IEF45p) as an independent check.
      On the imported null surface, compare the scalar distribution with the
      measure (IEF45e-1), including all marker IDs and invariant weights.
      Throughout the common future, compare every covariant stress component
      (IEF45g)--(IEF45i-3), and both the metric and inverse metric under the
      Jacobian pullback (IEF45j)--(IEF45l).  Later empirical-\(f\) comparison
      additionally requires the optional trajectory artifact described above.
      Require the forward and inverse maps to compose to the identity within
      the discretization error.  Equal coordinate timestamps or equal raw
      radial arrays are not a test of coordinate equivalence.
\item \textbf{Post-horizon EF continuation.}
      Continue only the EF evolution through \(a=0\), demonstrate active
      markers at \(r<r_H\), and pass the intrinsic accepted-step,
      post-horizon-duration, mass-accounting, and causal-excision tests
      (IEF57c)--(IEF57i).  This establishes a candidate only.  Vary the
      excision trigger and buffer independently, authenticate the repeats,
      and require (IEF57j) before calling the endpoint certified.  Outer-domain
      variation remains a separate convergence check and is not replaced by
      the excision certificate.
\end{enumerate}

The Schwarzschild calculation cannot reach \(a=0\) at finite Schwarzschild
coordinate time.  ``Agreement up to horizon formation'' therefore means that
the same-event, tensorially transformed inter-coordinate difference converges
to zero over the common exterior domain and through increasingly compact
pre-horizon slices.  Compactness milestones are useful for selecting slices,
but they are not a substitute for solving the event map.  A passing certificate
requires at least two radial/time resolutions and two independent quiet-PIC
realizations; missing map metadata, inverse closure, matter-support coverage,
full moments, or convergence evidence makes the result
\code{inconclusive}, never \code{pass}.

\subsubsection{Angular-momentum threshold protocol}

There is no universal critical angular momentum for neutral
Einstein--Vlasov collapse.  It depends on the rest-mass profile, radial
momentum distribution, angular momentum distribution, compactness, and on
which quantity is held fixed as angular momentum changes.  East's value
\(L_c/M\simeq0.67\) belongs to his charged (\(q_0=1.25\)), near-extremal family and
must not be used as the expected neutral VENUS answer.

\paragraph{One bound family and one fixed perturbation.}
The production scan begins with one validated static energy-cutoff shell from
Section~\ref{sec:ef-static-energy-cutoff-shell}.  The continuum parameters
\((C,L_0,y_c)\), target ADM mass, one even marker count, quiet-sampling
realization, and the inward-kick parameters \((\kappa,r_c)\) in (IEF-S26) are
fixed before the scan.  The kick is applied once and identically at every
scan point.  No Gaussian radial tail, radial-temperature retuning, or
point-dependent kick is permitted.

The resulting coordinate-neutral base artifact fixes every marker ID, source
ID, relative invariant rest weight, radius, and kicked radial momentum.
Introduce only the deterministic angular scale \(s_L>0\),
\[
 L_p(s_L)=s_LL_0,\qquad
 \eta_p(s_L)=\frac{s_L^2L_0^2}{2}.                 \tag{IEF59}
\]
For this delta-angular momentum shell, all markers have the same magnitude at
a fixed \(s_L\), while their isotropic angular directions still give zero net
spacetime angular momentum.  The driver interface is therefore
\codesource{Angular-family production interface}{\path{cpu/src/drivers/einstein_vlasov_wpic.cpp}}{Authenticated angular rescaling at fixed target ADM mass}
\begin{lstlisting}[basicstyle=\ttfamily\scriptsize]
--high-initial-condition energy_cutoff_shell
--high-energy-cutoff-ell0 ELL0
--high-angular-momentum-scale S_L
--high-target-initial-adm-mass M_TARGET
\end{lstlisting}
The bare \code{--high-specific-angular-momentum} replacement is deliberately
rejected for this IC, because it would erase the authenticated continuum
provenance.

Changing \(L\) at fixed rest weight changes particle energy and hence the
gravitating mass.  To isolate the centrifugal control, VENUS applies (IEF59),
then multiplies all invariant marker weights by one common positive
\(\alpha(s_L)\), leaves their phase space locations and relative weights
unchanged, solves the Schwarzschild radial constraint, and uses bracketed
bisection to impose
\[
 M_{\rm ADM}\!\left[\alpha(s_L),s_L\right]
       =M_{\rm ADM}^{\rm target}.                       \tag{IEF59a}
\]
The solve fails if the target cannot be reached on the exterior Schwarzschild
branch; it never clips the metric or accepts a supercritical initial SC slice.
Both \(\alpha\), the realized rest weight, the target and measured initial ADM
masses, and \(L/M_{\rm ADM}\) are written to the run summary.  Positive
binding is re-evaluated from the \emph{realized} fixed-ADM markers at every
\(s_L\): a continuum baseline that is bound at \(s_L=1,\kappa=0\) does not by
itself prove that every kicked/rescaled member is bound.

For an initially bound family the centrifugal barrier makes a virialized or
long-lived halo the expected sufficiently high-\(L\) outcome, whereas reducing
\(L\) favors collapse.  This is the physical ordering to test, not a label to
impose.  A scan that finds an unbound/dispersing high-\(L\) endpoint has not
measured the requested black hole--halo threshold and must be redesigned or
reported as such.

\paragraph{Preflight before long jobs.}
No final values of \((C,L_0,y_c,\kappa,r_c,M_{\rm ADM}^{\rm target})\),
\(R_B,R_E,T_{\rm donor}\), end time, or scan bracket are hard-coded by this
manual.  They are frozen only after local tests establish all of the
following:
\begin{enumerate}
\item the zero-kick \(s_L=1\) baseline reproduces the continuum sources,
      constraints, binding, current cancellation, and stable Lagrangian
      radii under radial/time/marker refinement;
\item the selected fixed kick remains positively bound and exterior on its
      initial SC slice, and a high-\(L\) endpoint shows bounded breathing and
      sustained Jeans/radius/profile stability rather than escape;
\item a lower-\(L\) endpoint forms an EF apparent horizon and continues for a
      resolved post-horizon interval without violating fixed-\(M\), \(b>0\),
      \(P>0\), or causal excision, thereby passing the intrinsic candidate
      gate; its independent excision-threshold and buffer repeats must then
      pass (IEF57j) before it is a certified threshold endpoint;
\item the chosen \(R_B,R_E,T_{\rm donor}\) pass every vacuum, TSC-clearance,
      crossing, eikonal, analytic-extension, and donor-history gate stated
      above for the largest production marker realization; and
\item short initialize-only and short dynamic runs produce complete,
      internally consistent IC, summary, restart, and map provenance before
      the cluster wall time is increased.
\end{enumerate}

First perform a coarse logarithmic or monotone scan to find one intrinsic
black hole candidate at low \(s_L\) and one genuinely virialized, positively
bound halo case at high \(s_L\).  The candidate may guide the scout but does
not yet close a scientific bracket.  Monotonicity must be checked from the
outcomes rather than assumed.  Only then bisect the same marker realization,
and certify the low endpoint before reporting the final bracket:
\begin{enumerate}
\item label a run \path{black_hole_candidate} only after the outer \(a=0\)
      crossing and every intrinsic condition (IEF57c)--(IEF57i) pass; label it
      \path{black_hole_certified} only after the authenticated independent
      sensitivity test (IEF57j) also passes;
\item classify a virialized halo only when the global family diagnostic has
      \(f_{\rm bind}>0\), bound Lagrangian radii and time-averaged profiles
      remain stable for several dynamical times; missing or nonpositive
      \(f_{\rm bind}\) fails this halo gate, while black hole classification
      remains independent of it;
\item distinguish a dispersing high-\(L\) cloud from a virialized halo using
      escaped rest weight and growing Lagrangian radii;
\item retain an ``undetermined'' class for runs whose near-critical dwell time
      exceeds the integration window;
\item repeat the terminal bracket at increased marker, radial, temporal, and
      domain resolution, and with the second authenticated quiet realization;
\item quote a converged bracket, not more digits than the resolution and PIC
      noise support.
\end{enumerate}
Near a type-I threshold the dwell time can increase strongly as the bracket
narrows.  Failure to form a horizon by a fixed short time is not, by itself,
evidence for a stable halo.

\paragraph{Which outputs decide the threshold.}
For a horizon-penetrating campaign, the threshold classifier consumes the EF
scan directory at each \(s_L\).  The associated SC calculation is the
physical, forward-time incoming solution used to build that point's null
initial data; it also provides pre-horizon comparison data.  Inside the
extractor, a separate scratch copy with reversed radial momenta is evolved
forward only to supply the negative-\(\tau\) part of the donor history, as
explained in (IEF45c-1).  The physical SC evolution is never replaced by that
scratch copy.  If a paired long SC run is retained, it may support the halo
diagnosis while \(a>0\), but the EF run alone decides post-horizon black hole
candidate evidence; the independent excision repeats decide whether that
candidate is certified.

An initial multi-\(s_L\) scout is only a search for a possible low-\(L\)
black hole/high-\(L\) halo bracket.  It is not a coordinate-equivalence
certificate, even though every EF point has an authenticated donor map.  The
separate \(2\times2\) parity workflow must already have issued the strict
certificate that authorizes the scout.  A scout bracket is likewise not a
final threshold until its black hole endpoint has the separate authenticated
excision sensitivity certificate and nearby points have been repeated at
finer resolution and with another quiet-PIC realization.  A bare
\path{black_hole_candidate} never closes the bracket.

For the required paired-coordinate cluster experiment, the Schwarzschild job
first writes a dense donor history and an authenticated solved-eikonal null
artifact.  The following is a symbolic template: every upper-case shell
variable is a value fixed by the successful local preflight, not a shipped
production default.
\codesource{SC donor and solved-null export interface}{\path{cpu/src/drivers/einstein_vlasov_wpic.cpp}}{Symbolic physical SC run and authenticated EF null-artifact export}
\begin{lstlisting}[basicstyle=\ttfamily\scriptsize]
--coordinates schwarzschild
--physical-ic common_ic.bin
--high-initial-condition energy_cutoff_shell
--high-energy-cutoff-amplitude "$STATIC_C"
--high-energy-cutoff-ell0 "$STATIC_ELL0"
--high-energy-cutoff-central-y "$STATIC_YC"
--high-energy-cutoff-power "$STATIC_K"
--high-radial-infall "$FIXED_KICK"
--high-infall-core-radius "$FIXED_KICK_CORE"
--high-angular-momentum-scale "$PARITY_S_L"
--high-target-initial-adm-mass "$TARGET_ADM"
--high-sampling-realization "$REALIZATION"
--export-ef-null-initial-data ingress_null.bin
--ef-null-donor-duration "$DONOR_DURATION"
--ef-map-worldtube-radius "$R_B"
--high-r-max "$R_E"
--high-radial-cells "$N_R"
\end{lstlisting}
The EF job must consume that artifact; giving it only the spacelike IC is a
hard error:
\codesource{EF solved-null import interface}{\path{cpu/src/drivers/einstein_vlasov_wpic.cpp}}{Symbolic EF production run consuming the authenticated null artifact}
\begin{lstlisting}[basicstyle=\ttfamily\scriptsize]
--coordinates eddington_finkelstein
--physical-ic common_ic.bin
--high-initial-condition energy_cutoff_shell
--high-energy-cutoff-amplitude "$STATIC_C"
--high-energy-cutoff-ell0 "$STATIC_ELL0"
--high-energy-cutoff-central-y "$STATIC_YC"
--high-energy-cutoff-power "$STATIC_K"
--high-radial-infall "$FIXED_KICK"
--high-infall-core-radius "$FIXED_KICK_CORE"
--high-angular-momentum-scale "$PARITY_S_L"
--high-target-initial-adm-mass "$TARGET_ADM"
--high-sampling-realization "$REALIZATION"
--ef-null-initial-data ingress_null.bin
--high-r-max "$R_E"
--high-radial-cells "$N_R"
\end{lstlisting}
The production parity directed acyclic graph (DAG), meaning the dependency
graph that orders extraction, paired runs, comparisons, and certification,
repeats this mapping for two resolutions and two
independent deterministic quiet-sampling realizations, builds convergence
evidence, and certifies the fine realization only after the same-event tests
above pass.  Only then is the angular scan released.  The coordinate-specific
realized marker hashes are expected to differ; equivalence is authenticated by
the common base IC plus the map-manifest graph, invariant marker data, and the
initial-null-slice empirical-measure audit and same-event tensor pullbacks.
After the shared exterior
comparison, only the EF run is continued into the trapped region.  The final
scientific result is a resolution- and realization-supported interval in
\(L/M_{\rm ADM}^{\rm target}\), together with the complete fixed-family
definition and outcome diagnostics, never a universal number.  Its low-\(L\)
endpoint must be \path{black_hole_certified}; an intrinsic candidate without
the authenticated threshold/buffer repeats is reported as inconclusive.
